\documentclass[11pt,letterpaper]{amsart}
\usepackage[T1]{fontenc}
\usepackage{lmodern}
\usepackage[margin=1.1in]{geometry}
\usepackage{amsmath,amssymb,amsthm,mathtools,mathrsfs}
\usepackage{microtype,booktabs,enumitem,needspace}
\usepackage{xcolor}
\usepackage{hyperref}
\hypersetup{colorlinks=true,linkcolor=blue!45!black,citecolor=blue!45!black,
 urlcolor=blue!45!black,pdftitle={Density of boundary entropy spectra for surface groups},
 pdfauthor={}}
\setlist[enumerate]{label=\textup{(\roman*)},leftmargin=2em}
\newtheorem{theorem}{Theorem}[section]
\newtheorem{lemma}[theorem]{Lemma}
\newtheorem{proposition}[theorem]{Proposition}
\newtheorem{corollary}[theorem]{Corollary}
\newtheorem{maintheorem}{Theorem}
\renewcommand{\themaintheorem}{\Alph{maintheorem}}
\theoremstyle{definition}
\newtheorem{definition}[theorem]{Definition}
\newtheorem{question}[theorem]{Question}
\theoremstyle{remark}
\newtheorem{remark}[theorem]{Remark}
\newtheorem{example}[theorem]{Example}
\numberwithin{equation}{section}
\DeclareMathOperator{\tr}{tr}
\DeclareMathOperator{\supp}{supp}
\DeclareMathOperator{\SL}{SL}
\DeclareMathOperator{\PSL}{PSL}
\DeclareMathOperator{\Hom}{Hom}
\newcommand{\CP}{\mathbb{CP}}
\newcommand{\RP}{\mathbb{RP}}
\newcommand{\HH}{\mathscr H_{\partial}}
\newcommand{\E}{\mathbb E}
\title[Boundary entropy spectra for surface groups]{Density of boundary entropy spectra\\for surface groups}
\date{2 November 2026}
\keywords{Poisson boundary, Furstenberg entropy, surface group, stationary measure,
Lyapunov exponent, entropy spectrum}
\begin{document}
\begin{abstract}
Let $\Gamma_g$ be a closed orientable surface group of genus $g\ge2$, and
let $\vartheta$ be a finitely supported probability measure whose support
generates $\Gamma_g$ as a semigroup. We prove that its boundary entropy
spectrum is dense in $[0,h_{\rm RW}(\Gamma_g,\vartheta)]$. The approximants
are essentially free projective boundaries of faithful representations
over the fixed field $\mathbb Q(i,\sqrt2)$; their matrix images may be
chosen dense in $\mathrm{SL}_2(\mathbb C)$, with stationary measures of
full support and exact dimension two. For the uniform eight-step walk
on the genus-two group, a comeager set of targets is attained exactly,
and the entire interval $[0,\frac12\log3]$ is attained. The proofs combine
an explicit rational representation curve, Galois approximation,
stationary-measure dimension formulas, and a category argument using
bifurcation currents. Applications include full realization for
coordinate products and continuum many nonisomorphic proper boundaries.
We also give a self-contained relative entropy inequality for joins,
a complete metric on the boundary lattice, and an explicit discontinuity
of representation entropy. Full realization for the original genus-two
walk remains open in this paper.
\end{abstract}
\maketitle

\section{Introduction}
\label{sec:introduction}

A Poisson boundary records the asymptotic information in a random walk.
Its measurable factors retain different amounts of that information,
measured by Furstenberg entropy. The boundary entropy realization problem
asks which amounts can occur. Even when the group and its increment law
are fixed, the distinction between approximating an entropy value and
realizing it exactly is substantial.

Let $G$ be a countable group and let $\vartheta$ be a probability measure
whose support generates $G$ as a semigroup. A $\vartheta$-boundary means a
measurable $G$-equivariant factor of the Poisson boundary of
$(G,\vartheta)$, equipped with its stationary probability measure. For a
stationary $G$-space $(B,\nu)$ we use the convention
\[
 h_\vartheta(B,\nu)
 =\sum_{g\in G}\vartheta(g)D(g_*\nu\Vert\nu),
 \qquad
 D(\alpha\Vert\beta)=\int\log\frac{d\alpha}{d\beta}\,d\alpha.
\]
This agrees, by change of variables, with
$-\sum_g\vartheta(g)\int\log(d(g^{-1}_*\nu)/d\nu)\,d\nu$.
Write
\[
 h_{\rm RW}(G,\vartheta)=\lim_{n\to\infty}\frac1nH(\vartheta^{*n}),
 \qquad
 \HH(G,\vartheta)=\{h_\vartheta(B,\nu):(B,\nu)\text{ is a }
 \vartheta\text{-boundary}\}.
\]
All logarithms are natural. For finite Shannon entropy, every boundary
entropy lies between zero and $h_{\rm RW}$, and the trivial and Poisson
boundaries attain the endpoints; see \cite{KV,Kaimanovich}.

Our principal result concerns the full closure of this spectrum, with
additional control of the boundaries realizing the approximating values.

\Needspace{13\baselineskip}
\begin{maintheorem}[Density for surface groups]\label{thm:all-surface-density}
Let $\Gamma_g$ be the fundamental group of a closed orientable surface of
genus $g\ge2$, and let $\vartheta$ be a finitely supported probability
measure whose support generates $\Gamma_g$ as a semigroup. Then
\[
 \overline{\HH(\Gamma_g,\vartheta)}
 =[0,h_{\rm RW}(\Gamma_g,\vartheta)].
\]
Every value in this interval can be approximated by the entropies of
essentially free projective boundaries associated to projectively
faithful representations
\[
 \Gamma_g\longrightarrow\mathrm{SL}_2\bigl(\mathbb Q(i,\sqrt2)\bigr).
\]
The representations may be chosen with dense image in
$\mathrm{SL}_2(\mathbb C)$, and their stationary measures on
$\CP^1$ may be chosen with full support and exact dimension two.
\end{maintheorem}

The measure $\vartheta$ remains fixed throughout this approximation.
The proof interpolates a Lyapunov exponent along a rational
representation curve connecting the trivial representation to a
cocompact Fuchsian representation. Galois conjugation supplies a dense
set of faithful parameters over one fixed number field. At these
parameters, the dimension theorem of Rapaport--Ren \cite{RR} identifies
boundary entropy with the smaller of $h_{\rm RW}$ and four times the
Lyapunov exponent. Finite-index restriction extends the construction
from genus two to every genus.

For exact realization we specialize to
\begin{equation}\label{eq:fixed-walk}
 \Gamma=\langle a_1,b_1,a_2,b_2\mid[a_1,b_1][a_2,b_2]=e\rangle,
 \qquad
 \mu=\frac18\sum_{s\in S}\delta_s,
 \quad S=\{a_1^{\pm1},b_1^{\pm1},a_2^{\pm1},b_2^{\pm1}\}.
\end{equation}
Set $h=h_{\rm RW}(\Gamma,\mu)$.

\Needspace{10\baselineskip}
\begin{maintheorem}[Generic and interval realization]\label{thm:overview-realization}
The spectrum $\HH(\Gamma,\mu)$ contains a dense $G_\delta$ subset of
$(0,h)$, all of whose values are attained by essentially free projective
boundaries. Moreover,
\begin{equation}\label{eq:main-interval}
 [0,\tfrac12\log3]\subseteq\HH(\Gamma,\mu).
\end{equation}
In particular $[0,c_*h]\subseteq\HH(\Gamma,\mu)$, where
$c_*=2\log3/(3\log7)=0.376383356\ldots$.
\end{maintheorem}

The generic statement follows from lower semicontinuity of entropy
and the support theorem for bifurcation currents of
Deroin--Dujardin \cite{DD}. A category argument selects points on a
comeager set of exact Lyapunov levels. The interval assertion then uses
a quotient onto $F_2\times F_2$ and the actual coordinate-update measure
induced by \eqref{eq:fixed-walk}. Those quotient boundaries have a
nontrivial kernel as $\Gamma$-spaces. Theorem~\ref{thm:generic-product-regularization}
shows more generally that two comeager boundary entropy spectra give
full realization for a coordinate product.

The normalization in Theorem~\ref{thm:overview-realization} follows from
\begin{equation}\label{eq:h-bounds}
 \tfrac12\log3\le h\le\tfrac34\log7.
\end{equation}
Indeed $F_4$ maps onto $\Gamma$ preserving the eight labeled increments,
whereas $(a_1,b_1,a_2,b_2)\mapsto(x,y,y,x)$ maps $\Gamma$ onto $F_2$
and sends $\mu$ to its simple random walk. Entropy decreases under
quotients. For the simple random walk on $F_r$, reduced length has speed
$1-1/r$, and the endpoint conditioned on length is uniform on its
sphere; hence its entropy is $(1-1/r)\log(2r-1)$.

A separate part of the paper concerns selection of limits in the lattice
of factors of one Poisson boundary. Its joins and meets are the generated
and intersected completed invariant $\sigma$-algebras, respectively.
We write $\mathsf h(A)=h_\mu(A)$ in the following statement.

\Needspace{10\baselineskip}
\begin{maintheorem}[Relative joins and entropy distance]\label{thm:overview-lattice}
For the fixed walk \eqref{eq:fixed-walk}, if $D$ is a common factor of
boundaries $A$ and $C$, then
\[
 \mathsf h(A\vee C)+\mathsf h(D)\le\mathsf h(A)+\mathsf h(C).
\]
The formula
\[
 d(A,C)=2\mathsf h(A\vee C)-\mathsf h(A)-\mathsf h(C)
\]
defines a complete metric on the boundary lattice modulo null sets.
Metric convergence implies strong convergence of conditional expectations.
If $\mathsf h(A_n)\to t$ and
$\sum_n[\mathsf h(A_n\vee A_{n+1})-\mathsf h(A_n)]<\infty$, then
\[
 \mathsf h\!\left(\bigwedge_N\bigvee_{n\ge N}A_n\right)=t.
\]
\end{maintheorem}

The proof of the relative join inequality uses the nonstationary entropy
functional and convolution identity of Gorodetski--Kleptsyn \cite{GK}.
We include the argument because the common-factor subtraction is useful
for the selection criterion. Completeness alone does not establish
closedness of the scalar entropy spectrum.

\subsection{Relation to previous work}
Dodds--Furman \cite[Theorem A]{DF} prove a nontrivial interval of boundary
entropies for surface groups. Their product theorem assumes genus at
least three \cite[Theorem B]{DF}; their continuity theorem imposes
$L^2$ conditions on a stationary density and its logarithm
\cite[Theorem 4.1]{DF}. The latter hypotheses do not cover the
circle-supported Fuchsian endpoint used here. Hochman--Solomyak
\cite{HS} and Rapaport--Ren \cite{RR} instead relate entropy to
stationary-measure dimension under arithmetic separation. We use these
formulas at dense algebraic parameters and interpolate only the Lyapunov
exponent, whose continuity follows from Bocker-Neto and Viana \cite{BV}.

The lattice of Poisson-boundary factors and entropy continuity along
monotone lattice limits are studied by Dodds \cite{Dodds}.
Tamuz--Zheng \cite[Section 4]{TZ} use ordinary join subadditivity in
their work on asymptotic entropy spectra. The relative inequality above
and its metric consequences are presented with self-contained proofs.
Brieussel--Oppelmayer \cite{BO} prove continuity for Bowen--Poisson
bundles; membership in that class of stationary actions does not itself
identify an action as a factor of the fixed Poisson boundary considered
here. We also distinguish monotone factor convergence from a general
compactness assertion; the relevant withdrawn claims in \cite{BT,BLT}
are discussed in Section~\ref{sec:remaining}.

\subsection{Scope and organization}
Theorem~\ref{thm:all-surface-density} excludes open entropy gaps for all
walks in its statement. Theorem~\ref{thm:overview-realization} leaves at
most a meager set of missing values for \eqref{eq:fixed-walk}; it does
not establish full realization or the interval $[0,h/2]$.
Theorem~\ref{thm:explicit-jump} gives a discontinuity of entropy on the
faithful non-elementary representation locus. The level-set results in
Sections~\ref{sec:level-test} and \ref{sec:saturation-test} explain why
parameter-space saturation alone does not settle every target.

Section~\ref{sec:curve} constructs the rational family, and
Section~\ref{sec:density} proves Theorem~\ref{thm:all-surface-density}.
Sections~\ref{sec:level-test}--\ref{sec:free-generic} establish generic
realization. The lattice results precede the applications in
Section~\ref{sec:applications-editorial}, which include
\eqref{eq:main-interval} and full realization for coordinate products.
The final sections give the discontinuity example, the remaining
levelwise obstructions, and precise open selection questions.

\section{A rational family and arithmetic faithfulness}\label{sec:curve}

For $z\in\mathbb C\setminus\{-1/2\}$ set
\[
 a(z)=1+2z,\qquad d(z)=\frac{1+2z^2}{1+2z},\qquad
 A(z)=\begin{pmatrix}a(z)&2z\\z&d(z)\end{pmatrix},\qquad
 B(z)=\begin{pmatrix}a(z)&-2z\\-z&d(z)\end{pmatrix}.
\]
Both matrices have determinant one. Define
\[
 K(z)=[A(z),B(z)],\qquad
 C(z)=(1+14z)I+zK(z),
\]
and the parameter domain
\[
 D=\{z\in\mathbb C:z\ne-1/2,\ \det C(z)\ne0\}.
\]
Since $\det C(z)$ is a rational function that is nonzero at $z=0$,
its zero set is finite. Thus $D$ is the complex plane minus finitely
many points and in particular is path connected.
For every $z\in D$ define a representation by
\begin{equation}\label{eq:explicit-curve}
 \rho_z(a_1)=A(z),\quad \rho_z(b_1)=B(z),\quad
 \rho_z(a_2)=C(z)B(z)C(z)^{-1},\quad
 \rho_z(b_2)=C(z)A(z)C(z)^{-1}.
\end{equation}
This is a rational family defined over $\mathbb Q$. In fact $C(z)$
commutes with $K(z)$, and hence the relator evaluates to
\[
 K(z)\,C(z)[B(z),A(z)]C(z)^{-1}
 =K(z)C(z)K(z)^{-1}C(z)^{-1}=I.
\]

\begin{lemma}[Trivial and Fuchsian endpoints]\label{lem:curve-endpoints}
Both $0,1$ belong to $D$. The representation $\rho_0$ is trivial,
and the projectivization of $\rho_1$ is a faithful cocompact Fuchsian
representation of the genus-two surface group.
\end{lemma}
\begin{proof}
At zero all four images equal $I$, and $C(0)=I$.
At one the exact matrix calculations are
\[
 A=\begin{pmatrix}3&2\\1&1\end{pmatrix},\quad
 B=\begin{pmatrix}3&-2\\-1&1\end{pmatrix},\quad
 K=\begin{pmatrix}-15&-56\\-4&-15\end{pmatrix},\quad
 C=\begin{pmatrix}0&-56\\-4&0\end{pmatrix}.
\]
In particular $\det C=-224\ne0$ and
\[
 (\operatorname{tr}A,\operatorname{tr}B,\operatorname{tr}AB)=(4,4,6),
 \qquad \operatorname{tr}[A,B]=-30<-2.
\]
The Fricke characterization of one-holed-torus holonomy
\cite[Theorem 4.1.2 and Section 4.4]{GoldmanFricke} says that this pair is the
holonomy of a hyperbolic one-holed torus with geodesic boundary.

Let $\ell$ be the axis of the hyperbolic element $K$.
The matrix $C=K-\tfrac12\operatorname{tr}(K)I$ has the same eigenlines
as $K$, opposite nonzero eigenvalues, and negative real determinant.
Its projective action, interpreted as an orientation reversing isometry
of $\mathbb H^2$, is reflection in $\ell$. To check the last assertion,
conjugate the two eigenlines to $0,\infty$: the projective matrix becomes
$\operatorname{diag}(1,-1)$, whose orientation reversing extension is
$w\mapsto-\overline w$, reflection in the imaginary axis.

Double the one-holed torus across its geodesic boundary. Its metric
extends smoothly across the boundary and gives a closed hyperbolic
surface of genus two. Developing the second half by reflection in
$\ell$ gives its two holonomy generators $CBC^{-1}$ and $CAC^{-1}$;
the reversal of their order makes the boundary commutator $K^{-1}$.
Consequently the holonomy of the double is exactly the projectivization
of \eqref{eq:explicit-curve} at $1$, with the prescribed presentation.
The holonomy of a closed hyperbolic surface is discrete and faithful.
\end{proof}

By openness of the quasifuchsian locus, there is an ordinary complex
neighborhood $V\subset D$ of $1$ such that $\rho_w$ is projectively
faithful for every $w\in V$ \cite[Section 4.3]{GoldmanComplex}.

\begin{lemma}[Galois approximation on the rational curve]\label{lem:curve-galois}
The set of parameters $z\in D\cap\mathbb Q(i,\sqrt2)$ such that
$\rho_z$ is projectively faithful and
\begin{equation}\label{eq:trace-nonreal}
 \operatorname{Im}\big((\operatorname{tr}A(z))^2\big)\ne0
\end{equation}
is dense in $D$.
\end{lemma}
\begin{proof}
Let $\sigma$ be the automorphism of $\mathbb Q(i,\sqrt2)$ fixing
$i$ and sending $\sqrt2$ to $-\sqrt2$. The map
\[
 \mathbb Q(i)^2\longrightarrow\mathbb C^2,\qquad
 (r,s)\longmapsto(r+s\sqrt2,r-s\sqrt2)
\]
has dense image, since it is the restriction of an invertible
complex-linear map and $\mathbb Q(i)$ is dense in $\mathbb C$.
Thus, for any nonempty open $U\subset D$, we can choose $z$ with
$z\in U$ and $\sigma(z)\in V$ simultaneously.
All matrix entries in \eqref{eq:explicit-curve} are rational functions
over $\mathbb Q$, so
\[
 \sigma(\rho_z(\gamma))=\rho_{\sigma(z)}(\gamma)
 \qquad(\gamma\in\Gamma).
\]
If $\rho_z(\gamma)=\pm I$, applying $\sigma$ gives the same equality
at $\sigma(z)$. Projective faithfulness there implies $\gamma=e$.
This proves projective faithfulness at $z$ and, in particular, ordinary
faithfulness into $\mathrm{SL}_2(\mathbb C)$.

Finally
\[
 \operatorname{tr}A(z)=\frac{2+4z+6z^2}{1+2z}
\]
is nonconstant holomorphic on $D$. The square of this function is
also nonconstant. By the open mapping theorem, its real locus has
empty interior. We may therefore first shrink $U$ to a nonempty
open subset satisfying \eqref{eq:trace-nonreal}, and then make the
same simultaneous approximation.
\end{proof}

For each parameter from Lemma~\ref{lem:curve-galois}, the image group
has no finite orbit on $\mathbb{CP}^1$: a subgroup with a finite orbit
is either finite or virtually solvable, whereas $\Gamma$ is neither.
Condition \eqref{eq:trace-nonreal} implies that $A(z)$ has eigenvalues
of distinct absolute values and thus is proximal. It also excludes an
invariant generalized circle. Indeed, after conjugating such a circle
to $\mathbb{RP}^1$, matrices preserving it have real or purely imaginary
traces, so all their squared traces are real. These checks verify
precisely strong irreducibility, proximality, and the circle exclusion
in Rapaport--Ren's theorem. Algebraic entries verify its Diophantine
assumption \cite[Lemma 6.1]{HS}; the same number-field norm argument
applies to complex entries, as recorded after Definition 1.1 in \cite{RR}.

\section{Density of the full entropy spectrum}\label{sec:density}

We first record the entropy identity used in the proof. Let $\pi$ be a
projectively faithful representation of a finitely supported measured
group $(G,\vartheta)$ into $\mathrm{SL}_2(\mathbb C)$, where
$\operatorname{supp}\vartheta$ generates $G$ as a semigroup. Suppose its image
is strongly irreducible, proximal, and preserves no generalized circle.
The unique stationary measure $\eta$ on $\CP^1$ defines a genuine
$\vartheta$-boundary: random translates contract to point masses, and
the limiting map factors through the Poisson boundary
\cite[Theorem 2.3]{DF}. This statement also applies to nondiscrete images.
Write $\chi$ for the top Lyapunov exponent of $\pi_*\vartheta$.

The Ledrappier--Young identity in \cite[Theorem 2.14]{RR} gives
\begin{equation}\label{eq:ledrappier-young}
 h_\vartheta(\CP^1,\eta)=2\chi\dim\eta.
\end{equation}
To identify its entropy convention, let $X_1$ be the first increment
and $L$ the limiting projective point. The quantity $H(p)-\Delta$ in
that theorem is
\[
 H(X_1)-H(X_1\mid L)=I(X_1;L)
 =\sum_g\vartheta(g)D(\pi(g)_*\eta\Vert\eta)
 =h_\vartheta(\CP^1,\eta).
\]
When the matrix entries are algebraic, the Diophantine property follows
from the number-field norm estimate; see the discussion following
\cite[Definition 1.1]{RR} and \cite[Lemma 6.1]{HS}.
Faithfulness preserves the atom weights of every convolution power,
so the matrix walk has asymptotic discrete entropy
$h_\vartheta=h_{\rm RW}(G,\vartheta)$. The dimension formula
\cite[Theorem 1.2]{RR} then reads
\begin{equation}\label{eq:general-good-entropy}
 \dim\eta=\min\{2,h_\vartheta/(2\chi)\},
 \qquad h_\vartheta(\CP^1,\eta)=\min\{h_\vartheta,4\chi\}.
\end{equation}
Rapaport--Ren use base-two logarithms; changing both entropy and
Lyapunov exponent to natural logarithms preserves these identities.

\begin{proof}[Proof of Theorem~\ref{thm:all-surface-density}: density and arithmetic control]
Put $K_0=\mathbb Q(i,\sqrt2)$ and
$h_\vartheta=h_{\rm RW}(\Gamma_g,\vartheta)$. Nonamenability of
$\Gamma_g$ and the entropy criterion imply $h_\vartheta>0$
\cite{KV}. First take $g=2$ and use the rational family $\rho_z$, $z\in D$,
from \eqref{eq:explicit-curve}. Let $\mathcal G\subset D\cap K_0$
be the dense set from Lemma~\ref{lem:curve-galois}. For every
$z\in\mathcal G$, the representation is projectively faithful and
its image is strongly irreducible, proximal, and preserves no generalized
circle. This set of parameters depends on the representations, not on
the probability weights or the support of $\vartheta$.

We record why these properties remain available in higher genus.
There is a subgroup $H_g\le\Gamma_2$ of index $g-1$ isomorphic to
$\Gamma_g$: for $g>2$, take the kernel of an epimorphism
$\Gamma_2\to\mathbb Z/(g-1)\mathbb Z$, and use the Euler characteristic
of its covering surface. For $g=2$, set $H_2=\Gamma_2$.
Fix an identification $\Gamma_g\cong H_g$ and set
$\pi_z=\rho_z|_{H_g}$. Each matrix entry of $\pi_z(\gamma)$ is a
rational function over $\mathbb Q$. The representation $\pi_0$ is
trivial, whereas $\pi_1$ is projectively faithful and cocompact
Fuchsian. For $z\in\mathcal G$, projective faithfulness and algebraic
entries are inherited by restriction.

The restricted image has no finite projective orbit: a subgroup of
$\mathrm{PSL}_2(\mathbb C)$ with a finite orbit is finite or virtually
solvable, whereas $H_g$ is neither. A positive power of the loxodromic
matrix $A(z)$ belongs to $\rho_z(H_g)$, so the restricted image is
proximal. To verify the circle condition, suppose that
$\rho_z(H_g)$ preserves a generalized circle $C$. The normal core
$N_g=\bigcap_{\gamma\in\Gamma_2}\gamma H_g\gamma^{-1}$ has finite
index in $\Gamma_2$ and is contained in $H_g$. Its image is again
non-elementary and contains loxodromic elements. Such a group can
preserve at most one generalized circle: the fixed points of its
loxodromic elements lie on every invariant circle, and non-elementarity
provides at least three distinct such points. Normality therefore
implies $\rho_z(\gamma)C=C$ for every $\gamma\in\Gamma_2$, since each
$\rho_z(\gamma)C$ is another $\rho_z(N_g)$-invariant circle.
This contradicts the circle exclusion for $\rho_z(\Gamma_2)$.
Thus the restricted image satisfies all three dynamical hypotheses
of Rapaport--Ren.

Now keep $g$ and $\vartheta$ fixed, and write
$\chi_\vartheta(z)$ for the top Lyapunov exponent of
$\pi_{z*}\vartheta$. Enumerate the finite support of $\vartheta$
as $\{s_1,\ldots,s_m\}$ and put $p_j=\vartheta(s_j)>0$.
The matrix tuple
$(\pi_z(s_1),\ldots,\pi_z(s_m))$ varies continuously in
$\mathrm{GL}(2,\mathbb C)^m$. Bocker-Neto and Viana's
Theorem~A~\cite{BV} proves continuity of both extremal Lyapunov
exponents for arbitrary complex two-dimensional matrix tuples with
positive probability weights. It applies at every point of $D$,
including reducible tuples and collisions of labeled matrix values.
Consequently $\chi_\vartheta$ is continuous on $D$, and
$\chi_\vartheta(0)=0$. No inversion symmetry is used here.

At $z=1$, the orbit of the cocompact Fuchsian group
$\pi_1(\Gamma_g)$ satisfies an upper counting bound $Ce^R$ in
$\mathbb H^2$. Thus, for every $q>1$,
\[
 Z_q=\sum_{\gamma\in\Gamma_g}
       e^{-q d(o,\pi_1(\gamma)o)}<\infty.
\]
The relative-entropy inequality against these normalized exponential
weights gives, for the $\vartheta$-random walk $X_n$,
\[
 H(\vartheta^{*n})
 \le q\,\mathbb E d(o,\pi_1(X_n)o)+\log Z_q.
\]
The displacement drift equals $2\chi_\vartheta(1)$. Dividing by $n$
and then letting $q\downarrow1$ yields
\[
 h_\vartheta\le2\chi_\vartheta(1).
\]
Since $D$ is path connected, every $t\in[0,h_\vartheta]$ therefore
has a parameter $z_t\in D$ with $4\chi_\vartheta(z_t)=t$.

For $z\in\mathcal G$, the support of $\pi_{z*}\vartheta$ generates
$\pi_z(\Gamma_g)$ as a semigroup. Its algebraic entries give the
Diophantine property, as recorded after Definition~1.1 of~\cite{RR}.
Faithfulness preserves every atom weight of every convolution power,
so its random-walk entropy is $h_\vartheta$. The stationary measure
$\eta_{z,\vartheta}$ on $\mathbb{CP}^1$ is a genuine
$\vartheta$-boundary by \cite[Theorem~2.3]{DF}; the equivalence in
\cite[Section~2.9]{RR} gives the required real Zariski density.
Applying Theorems~1.2 and 2.14 of~\cite{RR},
with the first-step conditional-entropy identification above,
gives
\[
 h_\vartheta(\mathbb{CP}^1,\eta_{z,\vartheta})
 =\min\{h_\vartheta,4\chi_\vartheta(z)\}.
\]
Approximate $z_t$ by parameters in $\mathcal G$ and use continuity
of $\chi_\vartheta$ to obtain entropies converging to $t$.
All boundary entropies are bounded by $h_\vartheta$, proving equality
of closures. Finally, $\eta_{z,\vartheta}$ is nonatomic, and each
nonidentity group element fixes at most two projective points.
Projective faithfulness and countability of $\Gamma_g$ imply essential
freeness.
\end{proof}

The following refinement completes the last assertion of
Theorem~\ref{thm:all-surface-density}.

\begin{corollary}\label{cor:dense-image-full-dimension}
The approximating boundaries constructed above
may all be chosen so that the representation has dense image in
$\mathrm{SL}_2(\mathbb C)$ and the stationary measure has full support
and exact dimension two on $\mathbb{CP}^1$.
\end{corollary}

\begin{proof}
Write $h_\vartheta=h_{\rm RW}(\Gamma_g,\vartheta)>0$.
It suffices to approximate targets $t\in(0,h_\vartheta)$, since the
endpoints are limits of interior targets. In the preceding construction,
choose the good parameters sufficiently close to $z_t$ that
$4\chi_\vartheta(z)<h_\vartheta$. Rapaport--Ren's dimension formula
then gives exact dimension two. If the projective image were discrete,
packing balls in $\mathbb H^3$ would give an orbit bound $Ce^{2R}$.
The preceding relative-entropy argument, now with $q>2$, would imply
$h_\vartheta\le4\chi_\vartheta(z)$, a contradiction.

Let $L$ be the closure of $\pi_z(\Gamma_g)$ in
$\mathrm{SL}_2(\mathbb C)$. Its Lie algebra $\mathfrak l$ is nonzero
because $L$ is not discrete. The condition that $\operatorname{Ad}(g)$
preserves the real vector subspace $\mathfrak l$ is a real algebraic
condition on $g$. Since $\pi_z(\Gamma_g)$ is real Zariski dense,
$\mathfrak l$ is invariant under all of
$\operatorname{Ad}(\mathrm{SL}_2(\mathbb C))$. It is therefore a
nonzero ideal in the simple real Lie algebra $\mathfrak{sl}_2(\mathbb C)$,
and must equal that Lie algebra. Hence $L$ is open; connectedness gives
$L=\mathrm{SL}_2(\mathbb C)$.

Finally, stationarity and semigroup generation make the closed support
of $\eta_{z,\vartheta}$ invariant under $\pi_z(\Gamma_g)$, and then
under its closure. Transitivity of $\mathrm{SL}_2(\mathbb C)$ on
$\mathbb{CP}^1$ implies that the support is the entire projective line.
\end{proof}

\begin{remark}
Exact dimension two does not imply absolute continuity with respect to
spherical area. No absolute-continuity assertion is needed for the
density theorem or its corollary.
\end{remark}
\subsection{The fixed genus-two walk}
For the rest of the paper, unless another measure is specified, use
$(\Gamma,\mu)$ from \eqref{eq:fixed-walk}. Let $\chi(z)$ denote the top
Lyapunov exponent of the eight labeled matrices $\rho_z(s)$, each with
weight $1/8$. Matrix collisions do not change this convention.
The proof above gives continuity of $\chi$ on $D$, with
$\chi(0)=0$ and $h\le2\chi(1)$. Write $Y_z=(\CP^1,\eta_z)$ for the
unique projective stationary boundary whenever the image is strongly
irreducible and proximal. At the dense algebraic faithful parameters
from Lemma~\ref{lem:curve-galois},
\begin{equation}\label{eq:entropy-at-good-parameters}
 h_\mu(Y_z)=\min\{h,4\chi(z)\}.
\end{equation}
This parameter set is countable. The connected parameter domain is
used to interpolate $\chi$, not to interpolate boundary entropy.

\begin{corollary}\label{thm:full-density}
For the fixed walk, $\overline{\HH(\Gamma,\mu)}=[0,h]$, with
approximating boundaries having all the properties in
Theorem~\ref{thm:all-surface-density}.
\end{corollary}
\begin{proof}
Apply Theorem~\ref{thm:all-surface-density} with $g=2$ and
$\vartheta=\mu$.
\end{proof}

\section{Geometry and separation of Lyapunov levels}
\label{sec:level-test}

Let
\[
 Q=\{z\in D:\Re z>0,\ \Im z>0\},\qquad
 N=\{z\in Q:\rho_z\text{ is not projectively faithful}\},
\]
where $D$ is the domain of the rational family. Throughout this section,
$h=h_{\rm RW}(\Gamma,\mu)$ and $\chi$ denotes its Lyapunov exponent.

\begin{proposition}[Faithful representations on exact Lyapunov levels]
\label{prop:exact-levels}
Every representation in $Q$ is non-elementary, proximal, and preserves
no generalized circle. The set $N$ is countable. For every
$0<c<\chi(1)$, the level
\[
 L_c=\{z\in Q:\chi(z)=c\}
\]
satisfies $\dim_H(L_c\setminus N)\ge1$. In particular, for every
$0<t<h$, the family contains a projectively faithful, real-Zariski-dense
representation with $\chi=t/4$.
\end{proposition}
\begin{proof}
For each nonidentity $\gamma\in\Gamma$, the equation
$\rho_z(\gamma)=\pm I$ has finitely many solutions in $D$:
its entries are rational functions, and neither identity holds
identically because $\rho_1$ is projectively faithful. Taking the
union over $\gamma$ proves countability of $N$.

Write
\[
 x(z)=\operatorname{tr}A(z)
 =\frac12+3z+\frac{3}{2(1+2z)}.
\]
For $z=u+iv\in Q$,
\[
 \Re x(z)=\frac12+3u+
 \frac{3(1+2u)}{2|1+2z|^2}>0,
 \qquad
 \Im x(z)=3v\frac{|1+2z|^2-1}{|1+2z|^2}>0.
\]
Thus $x(z)^2$ is nonreal. Consequently $A(z)$ is loxodromic,
and the image preserves no generalized circle. Moreover,
\[
 \operatorname{tr}[A(z),B(z)]
 =2-\frac{32z^4(z+2)^2}{(1+2z)^2}\ne2
 \qquad(z\in Q),
\]
so the two-generator reducibility criterion excludes a common fixed
point. If $A$ and $B$ preserved an unordered pair, each would fix
both points: an element interchanging the pair has trace zero,
whereas $\operatorname{tr}A=\operatorname{tr}B=x\ne0$.
This again contradicts irreducibility. A finite invariant set of at
least three points would give a finite projective image, contradicting
loxodromicity. Hence the family is non-elementary on $Q$.
Its real-Zariski density follows from \cite[Lemma~2.11]{RR}.

The set $Q$ is a quadrant with finitely many points removed and is
connected. Continuity of $\chi$ and the limits at $0$ and $1$ give
points $p_-,p_+\in Q$ with $\chi(p_-)<c<\chi(p_+)$.
A smooth embedded arc joining these points has a sufficiently narrow
tubular neighborhood, giving a bi-Lipschitz embedding
\[
 F:[0,1]\times[-\varepsilon,\varepsilon]\longrightarrow Q
\]
with $\chi(F(0,s))<c<\chi(F(1,s))$ for every $s$.
Each arc $F([0,1]\times\{s\})$ meets $L_c$. Therefore the
Lipschitz map obtained from $F^{-1}$ by projection onto its second
coordinate maps $L_c\cap\operatorname{im}F$ onto
$[-\varepsilon,\varepsilon]$. It follows that $\dim_H L_c\ge1$.
Removing the countable set $N$ preserves this lower bound.
Finally, $h\le2\chi(1)$ implies $t/4<\chi(1)$ for $0<t<h$.
\end{proof}

\paragraph{Open problem: separation on individual levels.}
Let $E\subset Q$ consist of the parameters for which $\rho_z(S)$
fails the weak Diophantine condition of \cite[Definition~1.1]{RR}.
Does
\begin{equation}\label{eq:open-level-separation}
 L_{t/4}\setminus(N\cup E)\ne\varnothing
 \qquad\text{for every }0<t<h?
\end{equation}
An affirmative answer, together with the dimension formula, would
prove full entropy realization. The estimate $\dim_H E<1$ would
suffice by Proposition~\ref{prop:exact-levels}. In contrast, density
or full planar measure of $Q\setminus E$ does not imply
\eqref{eq:open-level-separation}, since a prescribed level may lie
in an exceptional null set. The arithmetic construction supplies a
countable dense set of parameters outside $E$.

\subsection{Parametric separation and symmetric generating sets}

Solomyak--Takahashi \cite[Theorem~2.10]{ST} establishes a
zero-dimensional exceptional parameter set for real-analytic families
of contracting projective iterated function systems satisfying the
infinite-word nondegeneracy condition of their Definition~2.9.
Their Lemmas~2.13 and 2.15 obtain a uniform derivative order from
compactness of the analytic coding maps and then deduce the dimension
bound. Strict invariance of an open set supplies the required
contraction.

Symmetric matrix sets do not satisfy this strict-invariance hypothesis
\cite[Remark~1.3(1)]{ST}. In addition, the strong Diophantine condition
in that theorem separates distinct formal words, whereas symmetric
surface-group tuples have cancellations and the defining relation at
every parameter. Indexing matrices by distinct group elements removes
those permanent equalities but leaves the contraction and coding
hypotheses unverified. Rapaport--Ren \cite[Section~1.2]{RR} treats
parametric exceptional-set estimates as an open direction; its
pointwise theorem supplies no such estimate for this family.

\subsection{Kernel parameters and order of contact}

\begin{proposition}[Unbounded order of contact with the identity]
\label{prop:unbounded-contact}
If $z_0\in D$ is not projectively faithful, then for every integer
$k\ge0$ there is a nontrivial $w\in\Gamma$ such that
\[
 \rho_z(w)-I=O((z-z_0)^{k+1}),
 \qquad \rho_z(w)\not\equiv I.
\]
At $z_0=0$, iterated commutators of the first two standard generators
have this property.
\end{proposition}
\begin{proof}
The nontrivial normal subgroup $\ker\bar\rho_{z_0}$ contains a
free subgroup of rank two. Indeed, in a faithful Fuchsian realization,
conjugate a nontrivial kernel element by an element moving its fixed
point pair. Suitable powers of these independent hyperbolic elements
generate a free subgroup by ping-pong. Taking further squares gives
free generators $u,v$ with $\rho_{z_0}(u)=\rho_{z_0}(v)=I$.
Set $w_1=[u,v]$ and $w_{j+1}=[u,w_j]$. These words are nontrivial.
The elementary estimate
\[
 U(z)=I+O((z-z_0)^a),\quad V(z)=I+O((z-z_0)^b)
 \ \Longrightarrow\ [U(z),V(z)]=I+O((z-z_0)^{a+b})
\]
gives $\rho_z(w_j)-I=O((z-z_0)^{j+1})$.
Projective faithfulness at $1$ excludes an identically trivial word
matrix. At zero one can take $u=a_1,v=b_1$.
\end{proof}

\begin{proposition}[Density of kernel parameters]
\label{prop:dense-kernels}
The nonfaithful set $N$ is dense in
\[
 U_{\rm low}=\{z\in Q:\chi(z)<h/4\}.
\]
On no nonempty open subset of $U_{\rm low}$ can a fixed finite
number of derivatives separate every nontrivial word matrix from
the identity with positive lower bounds uniform in the parameter.
\end{proposition}
\begin{proof}
The family on $Q$ is non-elementary, generally faithful, and
nontrivial: countability of $N$ gives general faithfulness, and
$\operatorname{tr}^2A$ is nonconstant. The Sullivan--Bers stability
theorem, in the form of \cite[Theorem~2.6]{DD}, implies that any
open set consisting of faithful representations consists of discrete
faithful representations.

For a discrete faithful representation into $\mathrm{PSL}_2(\mathbb C)$,
its orbit in $\mathbb H^3$ is uniformly separated. Ball packing gives
\[
 \#\{\gamma:d(o,\rho(\gamma)o)\le R\}\le Ce^{2R}.
\]
The Gibbs entropy inequality with comparison weights proportional
to $e^{-q d(o,\rho(\gamma)o)}$, $q>2$, and the displacement drift
$2\chi$ imply $h\le2q\chi$. Letting $q\downarrow2$ gives
$h\le4\chi$. Thus no nonempty open subset of $U_{\rm low}$ can
consist entirely of faithful representations, proving density of $N$.

Given such an open subset and an integer $k$, choose $z_0\in N$
there. Proposition~\ref{prop:unbounded-contact} supplies a nontrivial
word for which
\[
 \partial_z^j(\rho_z(w)-I)\big|_{z=z_0}=0
 \qquad(0\le j\le k).
\]
This excludes even word-dependent positive lower bounds for the
maximum of these derivatives, uniform in the parameter. Restricting
the estimate to faithful parameters does not restore a positive
lower bound: faithful parameters approach $z_0$, and the derivatives
are continuous.
\end{proof}

\begin{corollary}[Nonharmonicity in the subcritical region]
\label{cor:no-flat-patches}
The function $\chi$ is harmonic on no nonempty open subset of
$U_{\rm low}$. In particular it is nowhere locally constant there.
\end{corollary}
\begin{proof}
Deroin--Dujardin \cite[Definition~3.2 and Theorem~3.9]{DD} identifies
the support of $dd^c\chi$ with the bifurcation locus. The admissibility
conditions in their Section~3 hold here: the family is holomorphic,
generally faithful, nontrivial, and everywhere non-elementary.
The fixed finitely supported generating measure satisfies their
exponential moment condition (8). The stability locus is the interior
of the discrete-faithful locus \cite[Theorem~2.6]{DD}, so the packing
inequality in the preceding proof excludes every stable point of
$U_{\rm low}$. Hence
$U_{\rm low}\subset\operatorname{supp}(dd^c\chi)$, which excludes
harmonicity on any nonempty open subset.
\end{proof}

The derivative obstruction concerns estimates uniform on complex open
sets. Equation~\eqref{eq:open-level-separation} asks instead for one
weakly Diophantine parameter on each prescribed level. Density of
kernel parameters in the ambient domain does not resolve that question.

\section{Generic exact realization}\label{sec:category}

Lower semicontinuity of entropy and the absence of constant Lyapunov
patches give exact realization for a comeager set of target values.

\begin{lemma}[Generic fibers of a continuous scalar map]
\label{lem:category-levels}
Let $X\subset\mathbb R^2$ be open, let $I$ be an open interval,
and let $p:X\to I$ be a continuous surjection which is constant
on no nonempty open subset. If $G\subset X$ is dense $G_\delta$,
there is a dense $G_\delta$ set $R\subset I$ such that, for each
$t\in R$, $G\cap p^{-1}(t)$ is a dense $G_\delta$ subset of the
nonempty perfect Polish space $p^{-1}(t)$. In particular this
intersection has cardinality continuum.
\end{lemma}
\begin{proof}
Let $E$ consist of the values attained at local extrema of $p$.
A countable base shows that $E$ is countable: every such value is
the supremum or infimum of $p$ on a basis element. Since every
level is closed and nowhere dense, $X'=X\setminus p^{-1}(E)$ is
dense $G_\delta$, and $I'=I\setminus E$ is dense $G_\delta$.
The restricted map $p:X'\to I'$ is a continuous open surjection.
Indeed a connected ball about a point of $X'$ contains values on
both sides of its value, and its image therefore contains an interval
about that value. Deleting $E$ preserves this neighborhood statement
in the relative topology.

Write $G\cap X'=\bigcap_nU_n$ with each $U_n$ open dense in $X'$,
and take a countable open base $(V_j)$ for $X'$. For each $n,j$,
$p(V_j\cap U_n)$ is open dense in the open set $p(V_j)$.
To check density, any nonempty relatively open set
$W\subset p(V_j)$ has a nonempty open inverse image in $V_j$,
which meets $U_n$. It follows that
$p(V_j)\setminus p(V_j\cap U_n)$ is nowhere dense in $I'$:
its closure has empty interior in $p(V_j)$ and can add only
points on the boundary of that open set. Remove the countable union
of these closures. The remaining set $R$ is dense $G_\delta$ in
$I'$ and hence in $I$. For $t\in R$, every $U_n$ meets every
nonempty basic open subset of $p^{-1}(t)$, so it is dense in that
fiber. The fiber is closed in the Polish space $X'$, and is therefore
Baire; this proves the dense $G_\delta$ conclusion.

Finally, the fiber has no isolated points. Otherwise a sufficiently
small punctured disk about an isolated fiber point would be connected
and contain no other point of the same level. The continuous function
$p-t$ would have constant sign there, making its center a local
extremum, contrary to $t\notin E$. A nonempty perfect Polish space
and each of its dense $G_\delta$ subsets have cardinality continuum.
\end{proof}


\begin{lemma}[Entropy saturation is residual in parameter space]
\label{lem:residual-saturation}
For the surface-group family, put
\[
 X=\{z\in Q:0<4\chi(z)<h\},\qquad
 e(z)=h_\mu(\CP^1,\eta_z),
\]
where $\eta_z$ is the unique stationary measure. There is a dense
$G_\delta$ subset $G\subset X$ on which $\rho_z$ is projectively
faithful and $e(z)=4\chi(z)$.
\end{lemma}
\begin{proof}
Non-elementarity throughout $Q$ gives weak continuity of $\eta_z$:
any weak limit along convergent parameters is stationary for the
limit tuple and hence equals its unique stationary measure.
Write $D_{\rm KL}$ for relative entropy of probability measures.
The change-of-variables form of Furstenberg entropy is
\[
 e(z)=\frac18\sum_{s\in S}
 D_{\rm KL}\bigl(\rho_z(s)_*\eta_z\,\|\,\eta_z\bigr).
\]
Each summand is lower semicontinuous in $z$. One can see this directly
from the variational identity, for probability measures $\alpha,\beta$
on the compact space $\CP^1$,
\[
 D_{\rm KL}(\alpha\|\beta)
 =\sup_{f\in C(\CP^1)}
 \left(\int f\,d\alpha-\log\int e^f\,d\beta\right).
\]
Thus $e$ is lower semicontinuous. The Ledrappier--Young identity
$e(z)=2\chi(z)\dim\eta_z$ gives $e(z)\le4\chi(z)$ without
any separation assumption. The dense faithful algebraic parameters
from Lemma~\ref{lem:curve-galois} satisfy equality on $X$, by
\eqref{eq:entropy-at-good-parameters}.

It follows that, for every $n\ge1$,
\[
 O_n=\{z\in X:e(z)>4\chi(z)-1/n\}
\]
is open and dense. Intersect these sets and remove the countable
nonfaithful set $N$. The result is a dense $G_\delta$ set $G$ with
the asserted properties. Notice that this argument does not assert
weak Diophantinity on $G$; it proves the entropy conclusion directly.
\end{proof}

\begin{theorem}[Generic exact realization throughout the full range]
\label{thm:generic-realization}
There is a dense $G_\delta$ set $R_\Gamma\subset(0,h)$ such that every
$t\in R_\Gamma$ is the entropy of an essentially free boundary of the
fixed $(\Gamma,\mu)$. Consequently
\[
 (0,h)\setminus\HH(\Gamma,\mu)\quad\text{is meager}.
\]
\end{theorem}
\begin{proof}
Use $X$ and $G$ from Lemma~\ref{lem:residual-saturation}, and set
$p(z)=4\chi(z)$. Proposition~\ref{prop:exact-levels} shows $p(X)=(0,h)$.
Corollary~\ref{cor:no-flat-patches}, based on the support theorem for
the bifurcation current, verifies the remaining hypothesis of
Lemma~\ref{lem:category-levels}. For $t$ in its dense $G_\delta$
set, select $z\in G$ with $p(z)=t$. Then $e(z)=t$.
The stationary projective space is a genuine boundary, as before.
Its measure is nonatomic, and projective faithfulness implies
essential freeness because a nonidentity M\"obius transformation
fixes at most two points.
\end{proof}

The exceptional target set is meager. This gives no assertion about its
Lebesgue measure, and a particular prescribed value can still
belong to it. A second construction will turn the corresponding
generic statement for a free group into a complete numerical interval
for the surface group.

\section{Generic realization for the free group}\label{sec:free-generic}

Write $F=\langle x,y\rangle$, let
\[
 \sigma=\tfrac14(\delta_x+\delta_{x^{-1}}+\delta_y+\delta_{y^{-1}}),
 \qquad h_2=h_{\mathrm{RW}}(F,\sigma)=\tfrac12\log3.
\]

\begin{proposition}\label{prop:free-generic}
There is a comeager set $R\subset(0,h_2)$ such that every $r\in R$
is the entropy of a genuine $\sigma$-boundary of $F$.  These boundaries
can be chosen to be essentially free projective boundaries of faithful
representations $F\longrightarrow\mathrm{SL}_2(\mathbb C)$.
\end{proposition}

\begin{proof}
Use the two matrices $A(z),B(z)$ already defined in the explicit curve,
now as the images of the two free generators.  Denote this representation
by $\pi_z$ and its top Lyapunov exponent for the four equally weighted
matrices by $\chi_F(z)$.  This is a separate exponent from the eight-step
surface-group exponent.

At $z=0$ the representation is trivial.  At $z=1$ it is the faithful
convex cocompact one-holed-torus representation checked above.
The faithful neighborhood $V$ for the surface-group family also gives
faithfulness for this free subgroup throughout $V$.  The same
$\mathbb Q(i,\sqrt2)$ Galois approximation proves that faithful
algebraic parameters are dense.  Every transcendental parameter is
faithful as well: a nontrivial word equal to $\pm I$ there would make
the corresponding rational-function identity hold identically,
contradicting projective faithfulness at $1$.

Use the same punctured first quadrant $Q$ as in
Proposition~\ref{prop:exact-levels}. Its trace and commutator calculations
involve only $A,B$, so they prove that $\pi_z$ is strongly irreducible,
proximal, and preserves no generalized circle throughout $Q$.

The function $\chi_F$ is continuous.  Its values approach zero as
$z\to0$ in $Q$, whereas $\chi_F(1)\ge h_2/2$ by the same
hyperbolic-plane orbit-counting argument used at the surface-group
endpoint.  Since a quadrant minus finitely many points is path
connected, every value in $(0,h_2/4)$ occurs.  Let
\[
 X=\{z\in Q:0<4\chi_F(z)<h_2\}.
\]
Then $X$ is open and the range of $4\chi_F$ on $X$ is $(0,h_2)$.

We check the absence of locally constant exponent patches, which is
essential for selection of entropy values rather than just parameters.
On every connected open subset of $X$, the family $\pi_z$ is an
admissible holomorphic family in Deroin--Dujardin's terminology: it
is nonelementary throughout, generally faithful, and nontrivial
because the squared trace $(\operatorname{tr}A(z))^2$ is nonconstant.  The finitely
supported generating measure $\sigma$ satisfies all their moment
conditions.  By \cite[Theorems 2.6 and 3.9]{DD}, the support of
$dd^c\chi_F$ is the bifurcation locus, and a stable open patch would
consist of discrete faithful representations.  For any such
representation, packing in $\mathbb H^3$ gives
\[
 h_2\le 4\chi_F.
\]
Indeed its discrete faithful orbit has at most $Ce^{2R}$ points in
an $R$-ball; the relative-entropy comparison with the weights
$e^{-q d(o,go)}$, $q>2$, gives $h_2\le2q\chi_F$, and then let
$q\downarrow2$.  This contradicts the definition of $X$.
Thus $\chi_F$ is not harmonic, and in particular not constant, on
any nonempty open subset of $X$.

For $z\in X$ let $e_F(z)$ be the Furstenberg entropy of the unique
stationary probability on $\mathbb{CP}^1$.  The stationary measures
vary weakly continuously.  The relative-entropy formula makes
$e_F$ lower semicontinuous, and the dimension formula gives
$e_F(z)\le4\chi_F(z)$.  At the dense faithful algebraic parameters,
Rapaport--Ren gives equality, since $4\chi_F<h_2$.
It follows that
\[
 G=\{z\in X:e_F(z)=4\chi_F(z)\}
\]
is a dense $G_\delta$ subset of $X$.  Intersecting it with the
comeager set of faithful parameters preserves this property.
Lemma~\ref{lem:category-levels}, applied to the continuous
nowhere locally constant function $4\chi_F:X\to(0,h_2)$ and this
comeager set, gives a comeager set $R$ of realized values.

The projective stationary spaces are genuine boundaries because
strong irreducibility and proximality give almost-sure convergence
of the random translates of their stationary measures to point
masses.  Their measures are nonatomic.  Faithfulness then makes
the $F$-actions essentially free, since each nonidentity projective
transformation fixes at most two points.
\end{proof}

\section{Boundary factors and entropy limits}\label{sec:remaining}

The results in this and the next two sections apply to any countable
group $G$ equipped with a finitely supported probability measure $\mu$
whose support generates $G$ as a semigroup. Let $(B,\nu)$ be its Poisson
boundary and put $H=h_\mu(B,\nu)$. We identify boundary factors with
complete $G$-invariant sub-$\sigma$-algebras of $B$, modulo null sets.
For such an algebra $A$, write $P_A$ for conditional expectation on
$L^2(B,\nu)$ and $h(A)$ for its Furstenberg entropy. Joins and meets
are taken in this lattice.

\subsection{Conditional derivatives and monotone limits}

For $g\in G$, put $r_g=d(g_*\nu)/d\nu$. If
$\mu^{*k}(g)>0$, stationarity gives
\[
 r_g\le \frac{1}{\mu^{*k}(g)}.
\]
An analogous bound for $g^{-1}$, together with the Radon--Nikodym
cocycle identity, gives a strictly positive lower bound for $r_g$.
Semigroup generation ensures that both bounds exist for every $g$.
In particular there are constants $0<c\le C<\infty$ such that
$c\le r_s\le C$ for every $s\in\operatorname{supp}\mu$.

The derivative on a factor pulls back to $P_A r_g$. With
$\phi(u)=u\log u$, the change-of-variables form of entropy is therefore
\begin{equation}\label{eq:conditional-entropy}
 h(A)=\sum_s\mu(s)\int_B\phi(P_A r_s)\,d\nu.
\end{equation}
Consequently, strong convergence $P_{A_n}\to P_A$ implies
$h(A_n)\to h(A)$, since $\phi$ is Lipschitz on $[c,C]$.
For the uniform eight-generator surface walk one can take
$c=1/8$ and $C=8$; symmetry also gives the equivalent formula
\[
 h(A)=-\frac18\sum_{s\in S}\int_B\log(P_A r_s)\,d\nu.
\]

Martingale convergence now proves continuity of entropy along
increasing joins and decreasing meets:
\[
 A_n\uparrow A\ \Longrightarrow\ h(A_n)\uparrow h(A),
 \qquad
 A_n\downarrow A\ \Longrightarrow\ h(A_n)\downarrow h(A).
\]
These are the monotone continuity statements in
\cite[Theorem C]{Dodds}. They concern factors of a fixed Poisson
boundary and do not assert entropy continuity for arbitrary convergent
representations. Theorem~\ref{thm:explicit-jump} exhibits precisely
that distinction.

\subsection{Compactness and scalar closedness}

Scalar density alone does not provide a convergent sequence of factors.
The compactness argument in \cite[Section 4, Theorem D]{Dodds} invokes
compactness of the space of all sub-$\sigma$-algebras in the strong
conditional-expectation topology, citing Beissner--T\"olle. That source
was withdrawn after a gap was identified in its compactness proof
\cite{BT}.

The general compactness assertion fails. On a standard product
probability space, let $X,Z_1,Z_2,\ldots$ be independent signs, with
$X$ fair and $\mathbb E Z_n=a\in(0,1)$. For
$\mathcal A_n=\sigma(XZ_n)$ one has
\[
 \mathbb E[X\mid\mathcal A_n]=aXZ_n,
 \qquad
 \bigl\|\mathbb E[X\mid\mathcal A_n]
       -\mathbb E[X\mid\mathcal A_m]\bigr\|_2^2
 =2a^2(1-a^2)\quad(m\ne n).
\]
Thus the corresponding conditional expectations have no strongly
convergent subsequence. This example concerns arbitrary subalgebras;
it does not disprove compactness for invariant subalgebras of a
particular Poisson boundary. It does show that the general compactness
argument cannot establish that specialized conclusion.

Burton--Lupini--Tamuz likewise withdrew the compactness claim in
Theorem 4.6 of the first version of their paper after finding a gap
in its proof \cite{BLT}. Their retained results on weak equivalence
and entropy continuity do not establish compactness of the boundary
entropy spectrum. None of these compactness claims is used below.

\section{Submodularity and a complete entropy metric}\label{sec:submodular}

Ordinary join-subadditivity appears in
\cite[Section 4, proof of Theorem 1.3]{TZ}. We give a relative form
with a common-factor correction. The proof uses nonstationary entropy;
the initial relatively independent joining need not be stationary,
and the limiting canonical joining need not be absolutely continuous
with respect to it.

\subsection{Nonstationary entropy and relative joins}

For a probability measure $m$ on a compact metrizable $G$-space define
\[
 \Phi_\mu(m)=\sum_g\mu(g)D(g_*m\Vert\mu*m).
\]
This is the nonstationary entropy of
\cite[Definition 2.3]{GK}. It is weakly lower semicontinuous,
satisfies $0\le\Phi_\mu(m)\le H(\mu)$, and agrees with Furstenberg
entropy when $m$ is stationary. The mixture identity
\[
 \sum_i p_iD(m_i\Vert q)
 =\sum_i p_iD(m_i\Vert\overline m)+D(\overline m\Vert q),
 \qquad \overline m=\sum_i p_i m_i,
\]
implies the nonstationary convolution identity
\begin{equation}\label{eq:Phi-cocycle}
 \Phi_{\mu*\kappa}(m)
 =\Phi_\mu(\kappa*m)+\Phi_\kappa(m).
\end{equation}
These identities are \cite[Lemma 2.8 and Proposition 2.10]{GK}.
Their proofs use only mixtures and invertibility, so apply to the
compact models used here. Indeed, condition on $g$ in the average of
$D(g_*h_*m\Vert\mu*\kappa*m)$ and use the intermediate mixture
$g_*(\kappa*m)$. The within-mixture term is $\Phi_\kappa(m)$ by
invariance of relative entropy under $g$; the remaining term is
$\Phi_\mu(\kappa*m)$. In particular,
\[
 \Phi_{\mu^{*n}}(m)
 =\sum_{k=0}^{n-1}\Phi_\mu(\mu^{*k}*m),
 \qquad h_{\mu^{*n}}(A)=n h(A).
\]

\begin{proposition}[Submodularity of boundary entropy]
\label{prop:boundary-submodularity}
If $A,C$ are $\mu$-boundaries and $D$ is a common boundary factor, then
\[
 h(A\vee C)+h(D)\le h(A)+h(C).
\]
In particular, this holds for $D=A\wedge C$.
\end{proposition}
\begin{proof}
Choose simultaneous compact metrizable models in which the actions and
the maps $A\to D$ and $C\to D$ are continuous. Such models can be
obtained by taking countable families of bounded Borel functions
separating points, closing these families under group translates and
pullbacks from $D$, and embedding the spaces by their coordinates in
countable products of compact intervals. The closures carry continuous
actions and factor maps; the measures are concentrated on the original
Borel images. Measurable actions may first be restricted to invariant
conull sets so that the group laws hold everywhere.

Disintegrate the stationary measures over $D$ and set
\[
 m_0=\int_D\nu^A_d\otimes\nu^C_d\,d\nu^D(d),
 \qquad m_n=\mu^{*n}*m_0.
\]
These are measures on the compact fiber product $A\times_D C$.
Writing $r_E(g,\cdot)=d(g_*\nu^E)/d\nu^E$ for $E=A,C,D$,
the conditional derivative in each fiber is the total-space derivative
divided by the base derivative. Hence
\[
 \frac{d(g_*m_0)}{dm_0}(a,c)
 =\frac{r_A(g,a)r_C(g,c)}{r_D(g,d)},
 \qquad d=\pi_A(a)=\pi_C(c).
\]
For each fixed word these derivatives are bounded above and away from
zero, as in Section~\ref{sec:remaining}. All logarithms in the finite
sums below are therefore integrable. Integrating against $g_*m_0$
and using its marginals gives
\[
 \sum_g\mu^{*n}(g)D(g_*m_0\Vert m_0)
 =n\bigl(h(A)+h(C)-h(D)\bigr).
\]
The mixture identity then yields
\begin{align*}
 \Phi_{\mu^{*n}}(m_0)
 &=n\bigl(h(A)+h(C)-h(D)\bigr)-D(m_n\Vert m_0)\\
 &\le n\bigl(h(A)+h(C)-h(D)\bigr).
\end{align*}

Along almost every random-walk path, the two marginals of the translate
of $m_0$ converge to point masses, since $A$ and $C$ are boundaries.
Every coupling of those marginals converges to the point mass on the
corresponding pair. Taking expectations shows that $m_n$ converges
weakly to the canonical boundary joining $\lambda$. Thus
\[
 h(A\vee C)=\Phi_\mu(\lambda)
 \le\liminf_{k\to\infty}\Phi_\mu(m_k).
\]
Cesaro averaging and \eqref{eq:Phi-cocycle} prove the assertion.
\end{proof}

\subsection{The entropy metric and summable perturbations}

\begin{proposition}[A complete metric on the boundary lattice]
\label{prop:boundary-entropy-metric}
On the boundary lattice modulo null sets,
\[
 d(A,C)=2h(A\vee C)-h(A)-h(C)
\]
is a complete metric. For every boundary $Z$,
\[
 d(A\vee Z,C\vee Z)\le d(A,C),
 \qquad |h(A)-h(C)|\le d(A,C).
\]
\end{proposition}
\begin{proof}
Monotonicity gives nonnegativity. If $d(A,C)=0$, then
$h(A\vee C)=h(A)=h(C)$, and strict entropy monotonicity between
boundary factors gives $A=C$. For completeness, strictness follows
from Lemma~\ref{lem:conditional-derivative-control} below: zero entropy
deficit between nested factors makes their projections agree on the
dense span of the derivatives, and hence on $L^2$. Submodularity over the common factor
$B$ gives
\[
 h(A\vee C)\le h(A\vee B\vee C)
 \le h(A\vee B)+h(B\vee C)-h(B),
\]
which is the triangle inequality. Applying submodularity over $A$
to $A\vee C$ and $A\vee Z$ gives
\[
 h(A\vee C\vee Z)-h(A\vee Z)
 \le h(A\vee C)-h(A).
\]
Add the analogous inequality with $C$ in place of $A$ to obtain
join contraction. The entropy estimate follows from
$h(A\vee C)\ge\max\{h(A),h(C)\}$.

Suppose first that $\sum_n d(A_n,A_{n+1})<\infty$. Put
\[
 T_N=\bigvee_{n\ge N}A_n,
 \qquad
 \epsilon_N=\sum_{n\ge N}
 \bigl(h(A_n\vee A_{n+1})-h(A_n)\bigr).
\]
These tails tend to zero. Submodularity bounds the entropy increment
from adding $A_{j+1}$ to $A_N\vee\cdots\vee A_j$ by
$h(A_j\vee A_{j+1})-h(A_j)$. Increasing-join continuity therefore gives
\begin{equation}\label{eq:tail-join-bound}
 0\le h(T_N)-h(A_N)\le\epsilon_N.
\end{equation}
Let $A=\bigwedge_N T_N$. Decreasing-meet continuity implies
$h(T_N)\to h(A)$, and $A\vee A_N\le T_N$ gives
\[
 d(A,A_N)
 \le h(T_N)-h(A)+h(T_N)-h(A_N)\longrightarrow0.
\]
Every Cauchy sequence has a subsequence with summable adjacent
distances. The preceding argument gives its limit, and the triangle
inequality gives the same limit for the original sequence.
\end{proof}

\begin{corollary}[Exact entropy from summable join costs]
\label{cor:summable-join-selection}
For any sequence of boundaries $(A_n)$, the following conditions
are equivalent:
\[
 \sum_n\bigl(h(A_n\vee A_{n+1})-h(A_n)\bigr)<\infty,
 \qquad
 \sum_n d(A_n,A_{n+1})<\infty.
\]
Under either condition, $A_n$ converges in $d$ to the genuine boundary
\[
 A=\bigwedge_{N=1}^{\infty}\bigvee_{n\ge N}A_n.
\]
In particular, if $h(A_n)\to t$, then $h(A)=t$.
\end{corollary}
\begin{proof}
Writing $c_n=h(A_n\vee A_{n+1})-h(A_n)$, we have $0\le c_n\le d(A_n,A_{n+1})$
and the telescoping identity
\[
 \sum_{n=1}^N d(A_n,A_{n+1})
 =2\sum_{n=1}^N c_n+h(A_1)-h(A_{N+1}).
\]
Nonnegativity of entropy proves equivalence. The tail-join argument
above proves convergence to the displayed $A$, and entropy is
$1$-Lipschitz for $d$.
\end{proof}

Completeness does not imply compactness or closedness of the scalar
entropy image. Further metric consequences are collected in
Corollary~\ref{cor:metric-applications}.

\section{Quantitative factor selection and relative extensions}
\label{sec:join-selection}

\subsection{Entropy deficit and conditional expectations}

The next estimate relates the entropy metric to the realization of
factors as conditional-expectation operators. It also gives explicit
constants for the eight-generator surface walk.

\begin{lemma}[Nested entropy deficit controls conditional derivatives]
\label{lem:conditional-derivative-control}
Suppose $A\subseteq C$, put $\Delta=h(C)-h(A)$, and choose $k\ge1$
such that $\mu^{*k}(g)>0$. Then
\[
 \|(P_C-P_A)r_g\|_2
 \le\frac{\sqrt{2k}}{\mu^{*k}(g)}\sqrt\Delta.
\]
For the uniform eight-generator walk, this implies
\[
 \|(P_C-P_A)r_g\|_2
 \le\sqrt{2\ell}\,8^\ell\sqrt\Delta
 \qquad(|g|=\ell\ge1).
\]
Moreover, the linear span of $\{r_g:g\in G\}$ is dense in
$L^2(B,\nu)$. For every $f\in L^2(B,\nu)$ there is consequently
a modulus $\omega_f(\delta)\to0$ as $\delta\downarrow0$, independent
of the nested pair, such that
\[
 \|(P_C-P_A)f\|_2\le\omega_f(\Delta).
\]
\end{lemma}
\begin{proof}
Set $p=\mu^{*k}(g)$. Since $r_g\le p^{-1}$, the function
$\phi(u)=u\log u$ satisfies $\phi''(u)\ge p$ on the range of
$u=P_A r_g$ and $v=P_C r_g$. Also $P_Av=u$, so the linear term
in Taylor's formula integrates to zero. Hence
\[
 \int\bigl(\phi(v)-\phi(u)\bigr)\,d\nu
 \ge\frac p2\|v-u\|_2^2.
\]
All the other terms in the entropy deficit for $\mu^{*k}$ are
nonnegative by conditional Jensen. The $g$ term alone therefore gives
$k\Delta\ge(p^2/2)\|v-u\|_2^2$. For the uniform surface walk,
$\mu^{*\ell}(g)\ge8^{-\ell}$.

For density of the derivative span, suppose $f$ is orthogonal to every
$r_g$. Let $\beta$ be the boundary map on the random-walk increment
space, and let $X_n$ be the product of the first $n$ increments.
The conditional law of $\beta$ given those increments is $(X_n)_*\nu$,
and thus
\[
 \mathbb E[f(\beta)\mid\omega_1,\ldots,\omega_n]
 =\int f r_{X_n}\,d\nu=0.
\]
Martingale convergence gives $f(\beta)=0$ and hence $f=0$.
Finally, approximate $f$ by a finite linear combination
$F=\sum_j b_jr_{g_j}$. The difference of nested orthogonal projections
has norm at most one, so the derivative estimates imply
\[
 \|(P_C-P_A)f\|_2
 \le\|f-F\|_2+
 \sum_j |b_j|\frac{\sqrt{2k_j}}{\mu^{*k_j}(g_j)}\sqrt\Delta,
\]
where $\mu^{*k_j}(g_j)>0$; constant terms contribute zero.
This proves the asserted uniform modulus.
\end{proof}

For arbitrary factors $A,C$, applying the lemma through $A\vee C$
and using $\sqrt{x}+\sqrt{y}\le\sqrt{2(x+y)}$ gives
\begin{equation}\label{eq:metric-derivative-control}
 \|(P_A-P_C)r_g\|_2
 \le\frac{2\sqrt{k}}{\mu^{*k}(g)}\sqrt{d(A,C)}.
\end{equation}
For the prescribed surface walk the right-hand constant is at most
$2\sqrt\ell\,8^\ell$. Density of the derivative span implies that
$d$-convergence entails strong convergence of conditional expectations.
No converse or continuity of the meet operation is asserted.

\begin{proposition}[Selection by controlled canonical joins]
Suppose $h(A_n)\to t$ and
\[
 \lim_{N\to\infty}\sup_{m,n\ge N}h(A_m\vee A_n)=t.
\]
Then $A_n$ converges in $d$ to a boundary $A$ with $h(A)=t$,
and $P_{A_n}\to P_A$ strongly. The same conclusions hold if
$\sum_n d(A_n,A_{n+1})<\infty$, equivalently if the one-sided
join costs in Corollary~\ref{cor:summable-join-selection} are summable.
\end{proposition}
\begin{proof}
The first hypothesis and $h(A_n)\to t$ imply
$\sup_{m,n\ge N}d(A_m,A_n)\to0$. Apply completeness and the
$1$-Lipschitz property of entropy, followed by
\eqref{eq:metric-derivative-control}. The summability assertion is
Corollary~\ref{cor:summable-join-selection}.
\end{proof}

For comparison, the operator estimate alone already proves selection
under the stronger condition $\sum_n\sqrt{d(A_n,A_{n+1})}<\infty$.
It makes the projections strongly Cauchy on the derivative span and
hence on $L^2$. A strong limit $P$ of conditional expectations is
self-adjoint, positive, unital and idempotent. If a bounded real $f$
lies in its range, Jensen gives $P(f^2)\ge f^2$; equality of integrals
gives equality. Polarization and bounded functional calculus show that
the bounded range is an algebra. It is dense in the full range, by
applying $P$ to bounded truncations, so $P$ is conditional expectation
onto a complete subalgebra. Composition $T_g$ by $g$ is bounded on
$L^2$ because $r_g$ is bounded; passing to the limit in
$(I-P_{A_n})T_gP_{A_n}=0$ proves invariance. Formula
\eqref{eq:conditional-entropy} then proves entropy convergence.
Submodularity replaces this square-root summability requirement by
the equivalent summable-cost conditions above.

\subsection{Relative extension and exact realization}

An alternative to a Cauchy selection is to extend factors by arbitrarily
small entropy increments. The relative entropy difference used below
is consistent with the fiber-entropy decomposition in
\cite[eqs.~(2.5)--(2.6)]{HT}; submodularity bounds the cost of adjoining
a factor by its entropy over the common factor.

\begin{proposition}[A relative extension criterion]\label{prop:relative-extension}
For factors $A,C$, put $r_A(C)=h(C)-h(A\wedge C)$.
Fix $T\in(0,H]$. If, for every $A$ with $h(A)<T$, there are
$C\not\le A$ with arbitrarily small $r_A(C)$, then
\[
 [0,T]\subseteq\HH(G,\mu).
\]
More precisely, if $t\in(0,H]$ is not realized, there is a factor
$A$ with $a=h(A)<t$ such that
\begin{equation}\label{eq:relative-obstruction}
 C\not\le A\quad\Longrightarrow\quad
 h(C)-h(A\wedge C)>t-a.
\end{equation}
In particular, $A$ contains every $C$ with $r_A(C)\le t-a$,
and therefore every $C$ with $h(C)\le t-a$.
\end{proposition}
\begin{proof}
Let $\mathcal P_t$ be the factors of entropy at most $t$. For any
nonempty chain $\mathcal C\subseteq\mathcal P_t$, put
$s=\sup_{E\in\mathcal C}h(E)$ and choose an increasing sequence
$E_n\in\mathcal C$ with $h(E_n)\to s$. Set $E=\bigvee_n E_n$.
For each $F\in\mathcal C$, comparability gives
\[
 h(F\vee E_n)=\max\{h(F),h(E_n)\}\longrightarrow s.
\]
Increasing-join continuity implies $h(F\vee E)=h(E)=s$, and
strict entropy monotonicity gives $F\le E$. Thus $E$ is an upper
bound in $\mathcal P_t$. The empty chain is bounded by the trivial
factor, so Zorn's lemma gives a maximal $A\in\mathcal P_t$.

If $t$ is not realized, then $a=h(A)<t$. For $C\not\le A$,
maximality implies $h(A\vee C)>t$. Submodularity yields
\[
 t<h(A\vee C)\le a+h(C)-h(A\wedge C),
\]
proving \eqref{eq:relative-obstruction}. Arbitrarily small relative
costs contradict this inequality for every unrealized $t\le T$.
The trivial boundary realizes zero.
\end{proof}

The same Zorn argument proves full realization if every proper factor
has strict extensions with arbitrarily small entropy increments.
Proposition~\ref{prop:relative-extension} gives a sufficient condition
for constructing such extensions through canonical joins. In particular,
taking $T=H/2$ reduces realization of the lower half-interval to small
relative costs outside every $A$ with $h(A)<H/2$. This allows candidates
of substantial absolute entropy when most of that entropy already
belongs to the common factor.

Neither scalar density nor the generic realization theorem verifies
this extension condition. The small-extension theorem of
Tamuz--Zheng \cite[Theorem 1.3]{TZ} applies to non-essentially-free
free-group boundaries and does not provide it for arbitrary faithful
surface-group boundaries. The selection and extension criteria above
are sufficient conditions for further exact realization, not additional
claims about the attained interval.

\section{Applications}\label{sec:applications-editorial}

We apply generic realization first to products and then to the prescribed
genus-two walk. The product theorem keeps the increment distribution
explicit: each step acts on one coordinate, so boundary entropy is a
weighted sum of the two coordinate entropies.

\subsection{Full realization for coordinate products}

\begin{theorem}[Exact realization for coordinate products]
\label{thm:generic-product-regularization}
For $i=1,2$, let $G_i$ be a countable group with a semigroup-generating
probability measure $\mu_i$ of finite Shannon entropy, and suppose
$0<h_i:=h_{\rm RW}(G_i,\mu_i)<\infty$. Assume that the boundary
entropy spectrum of $(G_i,\mu_i)$ contains a comeager subset
$R_i\subset(0,h_i)$. For $0<p<1$, define
\[
 \lambda_p=p(\mu_1\otimes\delta_e)
          +(1-p)(\delta_e\otimes\mu_2)
 \quad\text{on }G_1\times G_2.
\]
Then
\[
 h_{\rm RW}(G_1\times G_2,\lambda_p)=ph_1+(1-p)h_2,
 \qquad
 \mathscr H_\partial(G_1\times G_2,\lambda_p)
      =[0,ph_1+(1-p)h_2].
\]
If the values in $R_i$ are realized by essentially free boundaries,
then every value in the open interval is realized by an essentially
free boundary of $G_1\times G_2$.
\end{theorem}
\begin{proof}
Let $(B_i,\eta_i)$ be a $\mu_i$-boundary. The product measure
$\eta_1\otimes\eta_2$ is $\lambda_p$-stationary. Realize the walk
by independent coordinate choices, with probabilities $p,1-p$,
and independent increment streams of laws $\mu_1,\mu_2$.
Both coordinates are updated infinitely often almost surely. Along
the resulting sample path, the translates of each $\eta_i$ converge
to a point mass. Hence the translates of the product measure also
converge to a point mass. This proves that the stationary product
is a genuine boundary. The entropy formula is
\begin{equation}\label{eq:weighted-coordinate-product}
 h_{\lambda_p}(B_1\times B_2,\eta_1\otimes\eta_2)
 =p h_{\mu_1}(B_1,\eta_1)+(1-p)h_{\mu_2}(B_2,\eta_2),
\end{equation}
since the Radon--Nikodym derivative of a step is the derivative on
its active coordinate. The product action is essentially free if
both coordinate actions are essentially free.

For completeness, the random-walk entropy has the same weighted
formula. Let $N_n$ be the number of updates of the first coordinate
by time $n$, and let $Z_n$ be the position of the product walk.
Conditional on $N_n=k$, its two coordinates are independent with
laws $\mu_1^{*k}$ and $\mu_2^{*(n-k)}$. Therefore
\[
 H(Z_n\mid N_n)
 =\mathbb E\bigl[H(\mu_1^{*N_n})
                  +H(\mu_2^{*(n-N_n)})\bigr],
 \qquad
 0\le H(Z_n)-H(Z_n\mid N_n)\le\log(n+1).
\]
The law of large numbers, the definition of $h_i$, and the bounds
$H(\mu_i^{*k})\le kH(\mu_i)$ show that the conditional entropy
divided by $n$ converges to $ph_1+(1-p)h_2$. Thus this is the
random-walk entropy and bounds every boundary entropy from above.

Fix $0<t<ph_1+(1-p)h_2$. The interval
\[
 J_t=\left(\max\left\{0,\frac{t-(1-p)h_2}{p}\right\},
           \min\left\{h_1,\frac{t}{p}\right\}\right)
\]
is nonempty. On $J_t$, both $R_1$ and the inverse image of $R_2$
under $r\mapsto(t-pr)/(1-p)$ are comeager. Their intersection is
nonempty by the Baire category theorem. Selecting $r_1$ there and
setting $r_2=(t-pr_1)/(1-p)$ gives
$t=pr_1+(1-p)r_2$. Choose the corresponding boundaries and use
\eqref{eq:weighted-coordinate-product}. The trivial boundaries
realize zero, and the two Poisson boundaries realize the upper
endpoint.
\end{proof}

\begin{corollary}[Full realization on products of surface groups]
\label{cor:surface-product-full}
For the fixed $(\Gamma,\mu)$ and any $0<p<1$,
\[
 \mathscr H_\partial\!\left(\Gamma\times\Gamma,
       p(\mu\otimes\delta_e)+(1-p)(\delta_e\otimes\mu)\right)
 =[0,h].
\]
Every $t\in(0,h)$ is realized by an essentially free boundary.
Likewise, the coordinate product of two simple random walks on $F_2$
has boundary entropy spectrum $[0,\frac12\log3]$.
\end{corollary}
\begin{proof}
Apply Theorem~\ref{thm:generic-product-regularization} to
Theorem~\ref{thm:generic-realization} and
Proposition~\ref{prop:free-generic}, respectively.
\end{proof}

\subsection{An explicit interval for the genus-two walk}\label{sec:product-interval}

\begin{theorem}\label{thm:explicit-product-interval}
For the walk \eqref{eq:fixed-walk},
\[
 [0,\tfrac12\log3]\subseteq\HH(\Gamma,\mu).
\]
Consequently $[0,c_*h]\subseteq\HH(\Gamma,\mu)$, where
$c_*=2\log3/(3\log7)$. The boundaries in this interval construction
factor through $F_2\times F_2$ and are not essentially free as
$\Gamma$-spaces.
\end{theorem}
\begin{proof}
With $F=\langle x,y\rangle$ and its simple random walk $\sigma$, define
\[
 q:\Gamma\longrightarrow F\times F,\qquad
 a_1\mapsto(x,e),\quad a_2\mapsto(y,e),\quad
 b_1\mapsto(e,x),\quad b_2\mapsto(e,y).
\]
Both commutators in the defining relation map to the identity, and
$q$ is surjective. The push-forward measure is exactly
\[
 q_*\mu=\tfrac12(\sigma\otimes\delta_e)
              +\tfrac12(\delta_e\otimes\sigma).
\]
Proposition~\ref{prop:free-generic} and
Theorem~\ref{thm:generic-product-regularization} give every entropy in
$[0,h_2]$, where $h_2=\frac12\log3$. Pulling a boundary back along
$q$ preserves stationarity, the point-mass contraction property, and
entropy. In particular,
\begin{equation}\label{eq:coordinate-product-entropy}
 h_\mu(B_1\times B_2,\eta_1\otimes\eta_2)
 =\tfrac12h_\sigma(B_1,\eta_1)+\tfrac12h_\sigma(B_2,\eta_2).
\end{equation}
This proves the interval assertion for the prescribed measure.
The estimate $c_*h\le\frac12\log3$ follows from \eqref{eq:h-bounds}.
Finally $[a_1,b_1]$ is nontrivial in $\Gamma$ and lies in $\ker q$,
so all these pulled-back actions have a nontrivial kernel.
\end{proof}

The maximal entropy of this quotient walk is $h_2$. Thus this particular
construction cannot yield an entropy larger than $\frac12\log3$;
it imposes no upper bound on other boundary constructions for $\Gamma$.

\subsection{Abundance of proper boundaries}

\begin{corollary}[Abundance of proper, essentially free boundaries]
\label{cor:continuum-proper-boundaries}
Every nonempty open interval $J\subset(0,h)$ contains the entropy
values of continuum many pairwise nonisomorphic, essentially free
proper $\mu$-boundaries of $\Gamma$. They may all be realized on
$\mathbb{CP}^1$ with nonatomic, exact-dimensional stationary measures
of dimension $2$.
\end{corollary}
\begin{proof}
Let $R_\Gamma$ be the dense $G_\delta$ set in
Theorem~\ref{thm:generic-realization}. The set $R_\Gamma\cap J$ has
cardinality continuum. For each $t$ in this set, the theorem selects
a faithful projective representation with $4\chi=t$ and boundary
entropy $t$. The stationary measure is exact dimensional, and the
Ledrappier--Young identity
$h_\mu=2\chi\dim\eta$ gives $\dim\eta=2$
\cite[Theorem~2.14 and Appendix~A]{RR}.
Different values of $t$ give nonisomorphic measured $\Gamma$-spaces,
since Furstenberg entropy is preserved by equivariant
measure-preserving isomorphisms. Each boundary is proper because
its entropy is strictly smaller than $h$.
\end{proof}

The same conclusion holds for the simple random walk on $F_2$, with
$h$ replaced by $\frac12\log3$. The assertion concerns different
entropy values; it makes no claim that distinct representations on
one Lyapunov level yield nonisomorphic measured boundaries.

\subsection{Stable joins and endpoint convergence}

The entropy metric gives quantitative consequences for canonical factor
operations. In this subsection write $\mathsf h$ for entropy, let
$P$ be the full Poisson boundary, and let $\mathbf0$ be the trivial factor.

\begin{corollary}\label{cor:metric-applications}
The following hold for boundary factors of the fixed walk.
\begin{enumerate}
\item If $\sum_n\mathsf h(C_n)<\infty$, $J_N=\bigvee_{n\le N}C_n$,
and $J=\bigvee_{n\ge1}C_n$, then
\[
 d(J_N,J)\le\sum_{n>N}\mathsf h(C_n).
\]
\item Joins satisfy
\[
 d(A\vee C,A'\vee C')\le d(A,A')+d(C,C').
\]
\item $d(A,P)=h-\mathsf h(A)$ and $d(A,\mathbf0)=\mathsf h(A)$.
Thus entropy convergence to either endpoint implies convergence to that
endpoint in the entropy metric and strongly for conditional expectations.
\end{enumerate}
\end{corollary}
\begin{proof}
For (i), subadditivity bounds the entropy increase in every finite tail
by its total entropy; increasing-join continuity gives the same bound
for $J$. Since $J_N\le J$, their distance is the entropy difference.
For (ii), use the triangle inequality and contraction under joining a
fixed factor. The identities in (iii) follow from $A\vee P=P$ and
$A\vee\mathbf0=A$; the operator conclusion follows from
\eqref{eq:metric-derivative-control}.
\end{proof}

Corollary~\ref{cor:summable-join-selection} supplies a further application:
a scalar entropy approximation becomes exact whenever its successive
join costs are summable. The density theorem provides scalar
approximations, but does not by itself provide this additional control.

\section{An explicit failure of representation-entropy continuity}\label{sec:jump}

The same rational family gives an explicit obstruction to extending
entropy continuously across the real locus.

\begin{theorem}[Exact doubling in a convergent sequence]\label{thm:explicit-jump}
Let
\[
 z_* =\frac12-\frac7{20}\sqrt2,\qquad
 s_* =\frac12+\frac7{20}\sqrt2=1-z_*.
\]
For the fixed uniform measure \eqref{eq:fixed-walk}, $\rho_{z_*}$ is a faithful
algebraic representation into $\mathrm{SL}_2(\mathbb R)$. Write
$\chi_*:=\chi(\rho_{z_*})$. Then $\chi_*>0$, and there is a sequence
of faithful algebraic representations $\rho_{\zeta_n}$ converging to
$\rho_{z_*}$, with real-Zariski-dense images in
$\mathrm{SL}_2(\mathbb C)$, such that
\[
 h_\mu(Y_{z_*})=2\chi_*,\qquad
 \lim_{n\to\infty}h_\mu(Y_{\zeta_n})=4\chi_*.
\]
In particular, the boundary entropy is not continuous on the locus of
faithful non-elementary complex representations.
\end{theorem}

\begin{proof}
First, $7/5<\sqrt2<10/7$ gives $0<z_*<1/100$ and
$99/100<s_*<1$. For a real parameter $s\in(99/100,1)$, direct matrix
identities give
\[
 x(s):=\operatorname{tr}A(s)=\operatorname{tr}B(s)
       =2+\frac{6s^2}{1+2s}>2,
\]
\[
 Z(s):=\operatorname{tr}(A(s)B(s))
       =2+\frac{4s^2(s+2)^2}{(1+2s)^2}>2,
\]
and
\[
 \kappa(s):=\operatorname{tr}K(s)
       =2-\frac{32s^4(s+2)^2}{(1+2s)^2}<-28.
\]
For the last bound, $(s+2)/(1+2s)>1$ and
$32(99/100)^4>30$. The Fricke criterion \cite[Theorem 4.1.2]{GoldmanFricke} for a marked hyperbolic
one-holed torus applies to the positive trace triple
$(x(s),x(s),Z(s))$, since its commutator trace is below $-2$.
Consequently $\langle A(s),B(s)\rangle$ is the holonomy of a marked
one-holed torus with geodesic boundary corresponding to $K(s)$.

Writing $a=1+14s$ and using $\det K(s)=1$ gives
\[
 \det C(s)=a^2+as\kappa(s)+s^2
 <226-28\left(\frac{743}{50}\right)\left(\frac{99}{100}\right)<0.
\]
Thus the real projective transformation $C(s)$ is an
orientation-reversing isometry of $\mathbb H^2$ that centralizes the
hyperbolic element $K(s)$. It fixes the two endpoints of the boundary
axis individually. Such an isometry is the reflection across that axis
composed with a translation along it. The displayed representation
therefore doubles the one-holed torus across its geodesic boundary,
with an allowed twist. The resulting closed hyperbolic surface has
genus two, so $\rho_s$ is projectively discrete and faithful.
This is equivalently the standard Fuchsian combination theorem for
gluing two geodesic-boundary surfaces.

In particular $\rho_{s_*}$ is faithful. All entries of the family are
rational functions with rational coefficients. The automorphism
$\sqrt2\mapsto-\sqrt2$ sends $s_*$ to $z_*$. It preserves all matrix
relations, and therefore preserves faithfulness, including projective
faithfulness. It also shows that $C(z_*)$ is invertible. Thus
$\rho_{z_*}$ is faithful and algebraic. Its image is non-elementary:
faithfulness excludes a finite projective orbit, and
$\operatorname{tr}A(z_*)>2$ gives an unbounded element. Furstenberg's
positivity theorem gives $\chi_*>0$.

We next verify a uniform upper bound. For $0<z<1/100$,
\[
 \|A(z)-I\|_{\rm op}\le\|A(z)-I\|_{\rm F}
 \le\sqrt{13}\,z<\frac1{25}.
\]
Here the lower-right entry of $A-I$ is
$-2z(1-z)/(1+2z)$, whose absolute value is at most $2z$.
The same estimate applies to $B$. Since these matrices have determinant
one, their inverse operator norms equal their operator norms. Hence
\[
 \|A^{\pm1}\|,\|B^{\pm1}\|<\frac{26}{25},\qquad
 \|K\|<\left(\frac{26}{25}\right)^4<\frac65.
\]
Write $C=(1+14z)(I+uK)$ with $u=z/(1+14z)<1/100$. A Neumann-series
bound gives
\[
 \|C\|\,\|C^{-1}\|
 \le\frac{1+u\|K\|}{1-u\|K\|}<\frac{253}{247}.
\]
Each of the eight generator matrices consequently has norm below
\[
 \frac{26}{25}\frac{253}{247}<\frac{11}{10}.
\]
It follows by submultiplicativity that
$\chi_*\le\log(11/10)$.

The quotient
$a_1\mapsto x,b_1\mapsto y,a_2\mapsto y,b_2\mapsto x$
maps the fixed walk to the simple symmetric walk on $F_2$.
Therefore
\[
 h:=h_{\rm RW}(\Gamma,\mu)\ge\frac12\log3
 >4\log(11/10)\ge4\chi_*.
\]
The strict numerical comparison follows, without approximation, from
$11^8<3\cdot10^8$.

It remains to produce the approximating sequence. Let $V$ be a complex
neighborhood of $s_*$ in $D$ on which the family remains projectively
faithful, using openness of convex cocompact representations. For any
small complex open set $U$ near $z_*$, choose $u\in U$ and $v\in V$.
Approximate $(u+v)/2$ and $(u-v)/(2\sqrt2)$ by
$a,b\in\mathbb Q(i)$, respectively. Then
$\zeta=a+b\sqrt2\in U$ and its conjugate
$\sigma\zeta=a-b\sqrt2\in V$. Thus $\rho_\zeta$ is faithful and
algebraic. We can additionally arrange
$\operatorname{Im}(\operatorname{tr}A(\zeta))^2\ne0$, because this
is an open dense condition on the complex parameter domain.
Faithfulness excludes finite projective orbits; the nonreal squared
trace excludes compact image and invariant generalized circles.
The representations so obtained satisfy all dynamical hypotheses of
Rapaport--Ren, Theorem~1.2.

Choose $\zeta_n\to z_*$. Continuity of the Lyapunov exponent at the
non-elementary limit follows from weak convergence of the unique
stationary measures and Furstenberg's continuous norm-cocycle formula.
The algebraic Diophantine property and the dimension formulas give
\[
 h_\mu(Y_{z_*})=\min\{h,2\chi_*\}=2\chi_*,
\qquad
 h_\mu(Y_{\zeta_n})=\min\{h,4\chi(\rho_{\zeta_n})\}.
\]
The real formula uses \cite[Theorems 1.1 and 3.4]{HS}.
The complex formula uses \cite[Theorems 1.2 and 2.14]{RR},
together with the identity of Furstenberg entropy and the first-step
mutual information $H(p)-H(\omega_0\mid L)$. Since $4\chi_*<h$,
the second expression converges to $4\chi_*$, as claimed.
\end{proof}


The limiting representation preserves a generalized circle. Thus this
example disproves entropy continuity on the faithful non-elementary
locus; it does not assert a discontinuity with both the sequence and its
limit inside the real-Zariski-dense locus. All the stationary spaces here
are genuine boundary factors. Weak convergence of their stationary
measures follows from uniqueness at the non-elementary limit, so the
example directly separates weak convergence from entropy convergence.

\section{Levelwise saturation and holomorphic rigidity}
\label{sec:saturation-test}

\subsection{A limitation of semicontinuity and bifurcation support}

Full realization by the projective family would follow from finding,
for each $0<t<h$, a parameter with $4\chi=t$ and a stationary measure
of dimension $2$. The generic realization theorem establishes this
for a comeager set of targets. The following example shows that
lower semicontinuity, dense saturation, and full bifurcation support
alone do not imply saturation on every level.

Let $\Omega=\{r<|z|<R\}$, where $0<r<R<\sqrt h$, and put
$p(z)=|z|^2$. Given any countable dense set $A\subset\Omega$,
choose $c\in p(\Omega)\setminus p(A)$ and $0<\varepsilon<c$.
Define
\[
 F(z)=p(z)-\varepsilon\mathbf1_{\{p(z)=c\}}.
\]
Then $0\le F\le p$, $F$ is lower semicontinuous, and $p-F$ is
upper semicontinuous. Equality $F=p$ holds on $A$ and on an open
dense set of full planar measure. Moreover, $p$ is smooth and strictly
subharmonic, and $dd^cp$ has full support. Nevertheless, equality
fails everywhere on the level circle $\{p=c\}$, which has positive
logarithmic capacity. Thus these properties neither ensure saturation
on each level nor force the exceptional set to have zero capacity.
The example concerns the logical strength of the stated hypotheses;
it does not assert that this defect occurs for Furstenberg entropy.

\subsection{Holomorphic families with constant Lyapunov exponent}

\begin{proposition}[Rigidity on a subcritical constant-Lyapunov disk]
\label{prop:constant-lyapunov-disk}
Let $(\rho_w)_{w\in\Delta}$ be a holomorphic family of representations
of $\Gamma$ into $\mathrm{PSL}_2(\mathbb C)$, with $\rho_0$
faithful and non-elementary. If
\[
 \chi(\rho_w)=c<h/4\qquad(w\in\Delta),
\]
then, after shrinking $\Delta$, all representations in the family
are conjugate.
\end{proposition}
\begin{proof}
Shrink the disk so that every representation is non-elementary.
Faithfulness at zero implies general faithfulness: for each nontrivial
word, its identity locus is a proper analytic subset. If the family
were nontrivial up to conjugacy, it would be admissible for
\cite[Theorem~3.9]{DD}; finite support of $\mu$ supplies the
exponential moment condition. Constancy of $\chi$ makes $dd^c\chi$
vanish. The support theorem and \cite[Theorem~2.6]{DD} would then
make the family discrete and faithful. The packing inequality
$h\le4\chi=4c$ contradicts the hypothesis.
\end{proof}

The proposition applies to holomorphic curves in multiparameter
representation spaces as well. Thus a holomorphic deformation through
a faithful subcritical representation that stays on one Lyapunov level
can vary only by conjugacy, which preserves boundary entropy.
The existence and regularity of real paths within such levels remain
separate questions: continuity or H\"older continuity of $\chi$ alone
does not supply an analytic implicit-function theorem.

\section{Open realization questions}\label{sec:open-problems}

For the fixed genus-two walk, the established exact interval is
$[0,\frac12\log3]$, and the complement of the boundary entropy spectrum
in $(0,h)$ is meager. Two further selection statements would resolve
individual missing values.

\begin{question}[Saturation on each Lyapunov level]
For every $0<t<h$, does there exist a projectively faithful parameter
$z\in Q$ with
\[
 4\chi(z)=t,\qquad \dim\eta_z=2?
\]
\end{question}

An affirmative answer would give full realization by
\eqref{eq:ledrappier-young}. Proposition~\ref{prop:exact-levels} supplies
faithful points on every required level, while
Lemma~\ref{lem:residual-saturation} supplies dimension-two saturation
on a residual set of parameters. The missing assertion is their
intersection on every level. The example in Section~\ref{sec:saturation-test}
shows that lower semicontinuity, dense saturation, and full support of
the bifurcation current do not alone prove it. Dense kernel
specializations and unbounded word contact also prevent a uniform
finite-order transversality argument on the subcritical region.

The boundary lattice gives a different formulation. For a factor $A$,
set $r_A(C)=\mathsf h(C)-\mathsf h(A\wedge C)$, where $\mathsf h$
denotes boundary entropy.

\begin{question}[Arbitrarily small relative extensions]
Does every factor $A$ with $\mathsf h(A)<h/2$ satisfy
\[
 \inf_{C\not\le A}r_A(C)=0?
\]
Does the same assertion hold for every proper factor $A$?
\end{question}

Proposition~\ref{prop:relative-extension} proves that the first assertion
would give $[0,h/2]\subseteq\HH(\Gamma,\mu)$; the second would give
full realization. The common-factor term allows $C$ to have substantial
absolute entropy if most of it is already present in $A$. Neither the
density theorem nor completeness of the entropy metric currently
provides such an extension. The metric criterion instead asks for a
choice of approximants whose joins introduce a summable amount of
additional entropy. These are precise remaining requirements, rather
than consequences of scalar density.

\appendix
\section{An elementary continuity argument for symmetric walks}\label{sec:elementary-continuity}
For a finitely supported symmetric measure in $\mathrm{SL}_2(\mathbb C)$,
$\chi\ge0$ and
\[
 \chi=\inf_{n\ge1}\frac1n\sum_{s_1,\ldots,s_n}
 \mu(s_1)\cdots\mu(s_n)
 \log\|\rho(s_1\cdots s_n)\|.
\]
The right-hand side is an infimum of continuous functions of the
matrix tuple, proving upper semicontinuity. At a non-elementary
noncompact image, uniqueness of the stationary projective measure
implies its weak continuity: any subsequential limit remains
stationary for the limiting tuple. The Furstenberg integral formula,
whose integrand is jointly continuous on the compact projective space,
then gives continuity of $\chi$ at that tuple.
At a relatively compact image, $\chi=0$. At an image fixing one
projective point, conjugate to upper triangular matrices. The two
diagonal exponents are the averages of opposite log-modulus
characters; inversion symmetry makes both zero. The Lyapunov
exponents of a triangular cocycle are its diagonal exponents, so
$\chi=0$. If precisely a two-point set is preserved, conjugate to
monomial matrices. If the two points are not interchanged, the same
argument applies. Otherwise the sign process on the two coordinate
axes is a finite irreducible Markov chain with uniform stationary
law; the two opposite logarithmic increments have zero averaged
drift, again giving $\chi=0$. A finite orbit of size at least three
forces a finite projective image and hence $\chi=0$.
Upper semicontinuity and nonnegativity now give continuity at all
remaining tuples. This proof uses no entropy regularity assertion.

\begin{thebibliography}{99}

\bibitem{BT}
P. Beissner and J. M. T\"olle,
\emph{A compact topology for $\sigma$-algebra convergence},
preprint, 2018, \href{https://arxiv.org/abs/1802.05920}{arXiv:1802.05920}.
Withdrawn in version~3, 19 May 2021, because of a gap in the proof of
Theorem~2.6, Claim~3.

\bibitem{BV}
C. Bocker-Neto and M. Viana,
\emph{Continuity of Lyapunov exponents for random two-dimensional matrices},
Ergodic Theory Dynam. Systems \textbf{37} (2017), no.~5, 1413--1442.
\href{https://doi.org/10.1017/etds.2015.116}{doi:10.1017/etds.2015.116}.

\bibitem{BO}
J. Brieussel and H. Oppelmayer,
\emph{Continuity of Furstenberg entropy on Bowen Poisson bundles},
preprint, 2026,
\href{https://arxiv.org/abs/2610.29561v1}{arXiv:2610.29561v1}.

\bibitem{BLT}
P. Burton, M. Lupini, and O. Tamuz,
\emph{Weak equivalence of stationary actions and the entropy realization problem},
preprint, 2016,
\href{https://arxiv.org/abs/1603.05013v2}{arXiv:1603.05013v2}.
Version~2, 2 November 2016, withdraws Theorem~4.6 of version~1.

\bibitem{DD}
B. Deroin and R. Dujardin,
\emph{Random walks, Kleinian groups, and bifurcation currents},
Invent. Math. \textbf{190} (2012), no.~1, 57--118.
\href{https://doi.org/10.1007/s00222-012-0376-5}{doi:10.1007/s00222-012-0376-5}.

\bibitem{Dodds}
S. Dodds,
\emph{On the lattice of boundaries and the entropy spectrum of hyperbolic groups},
preprint, 2022,
\href{https://arxiv.org/abs/2211.04554v3}{arXiv:2211.04554v3}.

\bibitem{DF}
S. Dodds and A. Furman,
\emph{Quotients of Poisson boundaries, entropy, and spectral gap},
preprint, 2025,
\href{https://arxiv.org/abs/2504.09926v1}{arXiv:2504.09926v1}.

\bibitem{GoldmanComplex}
W. M. Goldman,
\emph{The complex-symplectic geometry of $\mathrm{SL}(2,\mathbb C)$-characters
over surfaces},
in \emph{Algebraic groups and arithmetic}, Tata Inst. Fund. Res.,
Mumbai, 2004, pp.~375--407.
\href{https://arxiv.org/abs/math/0304307v2}{arXiv:math/0304307v2}.

\bibitem{GoldmanFricke}
W. M. Goldman,
\emph{Trace coordinates on Fricke spaces of some simple hyperbolic surfaces},
in \emph{Handbook of Teichm\"uller theory}, Vol.~II,
IRMA Lect. Math. Theor. Phys., vol.~13, Eur. Math. Soc.,
Z\"urich, 2009, pp.~611--684.
\href{https://doi.org/10.4171/055-1/16}{doi:10.4171/055-1/16}.

\bibitem{GK}
A. Gorodetski and V. Kleptsyn,
\emph{Non-stationary version of Furstenberg's theorem on random matrix products},
J. Assoc. Math. Res. \textbf{4} (2026), no.~1, 1--41.
\href{https://doi.org/10.56994/JAMR.004.001.001}{doi:10.56994/JAMR.004.001.001}.

\bibitem{HT}
Y. Hartman and O. Tamuz,
\emph{Furstenberg entropy realizations for virtually free groups and lamplighter groups},
J. Anal. Math. \textbf{126} (2015), no.~1, 227--257.
\href{https://doi.org/10.1007/s11854-015-0016-2}{doi:10.1007/s11854-015-0016-2}.

\bibitem{HS}
M. Hochman and B. Solomyak,
\emph{On the dimension of Furstenberg measure for
$\mathrm{SL}_2(\mathbb R)$ random matrix products},
Invent. Math. \textbf{210} (2017), no.~3, 815--875.
\href{https://doi.org/10.1007/s00222-017-0740-6}{doi:10.1007/s00222-017-0740-6}.

\bibitem{Kaimanovich}
V. A. Kaimanovich,
\emph{The Poisson formula for groups with hyperbolic properties},
Ann. of Math. (2) \textbf{152} (2000), no.~3, 659--692.
\href{https://doi.org/10.2307/2661351}{doi:10.2307/2661351}.

\bibitem{KV}
V. A. Kaimanovich and A. M. Vershik,
\emph{Random walks on discrete groups: boundary and entropy},
Ann. Probab. \textbf{11} (1983), no.~3, 457--490.
\href{https://doi.org/10.1214/aop/1176993497}{doi:10.1214/aop/1176993497}.

\bibitem{RR}
A. Rapaport and H. Ren,
\emph{Dimension of Furstenberg measures on $\mathbb{CP}^1$},
preprint, 2025,
\href{https://arxiv.org/abs/2511.00729v2}{arXiv:2511.00729v2}.

\bibitem{ST}
B. Solomyak and Y. Takahashi,
\emph{Diophantine property of matrices and attractors of projective
iterated function systems in $\mathbb{RP}^1$},
Int. Math. Res. Not. IMRN \textbf{2021}, no.~16, 12639--12669.
\href{https://doi.org/10.1093/imrn/rnz309}{doi:10.1093/imrn/rnz309}.

\bibitem{TZ}
O. Tamuz and T. Zheng,
\emph{On the spectrum of asymptotic entropies of random walks},
Groups Geom. Dyn. \textbf{19} (2025), no.~3, 879--897.
\href{https://doi.org/10.4171/GGD/798}{doi:10.4171/GGD/798}.

\end{thebibliography}

\end{document}

